\documentclass[conference]{IEEEtran}

% T1 fontenc: gives Times (ptm) real small-caps coverage for the glyphs that
% otherwise fall back to TS1 inside IEEEtran's small-caps table captions
% (source of the "Font shape TS1/ptm/m/sc undefined" warning).
\usepackage[T1]{fontenc}
\usepackage{cite}
\usepackage{amsmath,amssymb,amsfonts}
\usepackage{algorithmic}
\usepackage{graphicx}
\usepackage{textcomp}
\usepackage{xcolor}
\usepackage{booktabs}
\usepackage{url}
\usepackage{array}
\usepackage{tikz}
\usetikzlibrary{arrows.meta,positioning,fit,backgrounds,calc,shapes.geometric,shapes.misc}

\begin{document}

\title{Cluster-Routed Hard-Negative Mining for UAV Cross-View Visual-Localization}

\author{
\IEEEauthorblockN{Kim-Phuong Phung}
\IEEEauthorblockA{
% TODO: department/faculty name and email
\textit{Le Quy Don Technical University}\\
Hanoi, Vietnam\\
email@example.com}
\and
\IEEEauthorblockN{Quang-Uy Nguyen}
\IEEEauthorblockA{
% TODO: department/faculty name and email
\textit{Le Quy Don Technical University}\\
Hanoi, Vietnam\\
email@example.com}
}

\maketitle

\begin{abstract}
UAV visual localization requires retrieving the correct geographic position
of an aerial query image from a large georeferenced satellite map. This paper
presents a cluster-routed descriptor pipeline for scalable UAV
cross-view localization. The system builds on pretrained transformer and ConvNeXt
backbones whose raw, L2-normalized output embedding serves directly as the
retrieval descriptor — used for descriptor matching, FAISS K-means routing,
and candidate scoring alike, with no projection, clustering, or routing head;
the descriptor is adapted by fine-tuning the last backbone stages. The same
cluster structure that routes queries at evaluation time also organizes
training: hard negatives are mined from a query's own cluster and its nearest
clusters, with co-positive exclusion preventing false negatives, so the
contrastive signal concentrates on the visually repetitive candidates where
retrieval errors actually occur. The loss combines paired in-batch InfoNCE,
descriptor variance regularization, same-cluster hard-negative suppression,
and soft cluster consistency. During
evaluation, the Game4Loc pre-tiled satellite gallery (256-pixel tiles at
multiple zoom levels) is clustered in the retrieval embedding space
and searched through nearest-cluster candidate selection, with an optional
vector-quantized semantic re-ranking of the top candidates. The pipeline
includes checkpointed
centroids, cached query and tile embeddings, cluster visualization, and
query-level inspection for reproducible training and evaluation.
\end{abstract}

\begin{IEEEkeywords}
UAV localization, cross-view retrieval, visual place recognition, metric
learning, FAISS, Swin Transformer, clustering.
\end{IEEEkeywords}

\section{Introduction}
Unmanned aerial vehicle (UAV) localization from imagery is a fundamental
capability for autonomous navigation, inspection, search, and mapping.
Operational systems must infer the position of a UAV image by matching it
against a geographically referenced satellite or orthomosaic map. This is
commonly formulated as cross-view retrieval and metric learning: the query image is
captured from an aerial perspective, while the database imagery is a set of
top-down map crops. The goal is to learn an embedding in which geographically
corresponding views are close and unrelated map locations are far apart, as in
deep metric-learning systems based on contrastive, triplet, and
classification-style objectives \cite{hadsell2006dimensionality,schroff2015facenet,sohn2016improved}.

This task remains challenging due to the substantial appearance gap between
UAV and satellite views, the large number of candidate map locations, and the
presence of many visually repetitive regions. A flat descriptor can support fast
global retrieval but may not provide sufficient local discrimination inside
visually similar areas. A fully specialized descriptor, however, is difficult
to apply efficiently over a large map.

Recent work has shown that transformer backbones are effective for UAV and
cross-view localization because self-attention can aggregate long-range layout
cues across roads, buildings, fields, and water boundaries. However, the
retrieval errors that remain are concentrated among visually repetitive
candidates: nearby tiles sharing similar road grids, rooftops, vegetation, or
water boundaries that a flat in-batch contrastive objective rarely contrasts
against one another. Randomly drawn or in-batch negatives are overwhelmingly
easy, so the descriptor receives little gradient signal exactly where
discrimination matters most — inside the neighborhoods of look-alike
locations.

The proposed approach formulates localization as learned routing followed by
candidate scoring in the same embedding space. The backbone's L2-normalized
output embedding is trained — by unfreezing its last stages — to place
corresponding anchor-positive crops into the
same K-means cluster \cite{macqueen1967methods}. At evaluation time, this
descriptor first routes a query to nearby map clusters and then scores
candidate tiles with the same descriptor. The central training mechanism is
that the same cluster structure also organizes hard-negative mining: negatives
are drawn from a query's own cluster and its nearest clusters, ranked by
descriptor similarity, with co-positive exclusion preventing samples that
share a positive tile from being pushed apart as false negatives. The
contrastive signal therefore concentrates on exactly the look-alike
neighborhoods where retrieval fails. This design keeps map-side computation
cacheable and avoids training separate routing and scoring heads.

The contributions are:
\begin{itemize}
    \item Cluster-structured hard-negative mining: the K-means partition
    learned for routing doubles as the negative-sampling structure — hard
    negatives come from a query's own cluster and its nearest clusters,
    ranked by descriptor similarity, with co-positive exclusion to prevent
    false negatives — focusing the contrastive objective on visually
    repetitive look-alike regions.
    \item A FAISS-based clustered coarse-to-fine search strategy for fast
    large-map localization, where Game4Loc pre-tiled gallery images are
    embedded, clustered, and searched through nearest-cluster candidate
    selection before descriptor scoring, with an optional vector-quantized
    semantic re-ranking stage for top-candidate ordering.
    \item A full training and evaluation pipeline for comparing backbones,
    scoring modes, clustering settings, and reference approaches, with cached
    map/query embeddings, checkpointed centroids, evaluation
    history, and query-level inspection.
\end{itemize}

\section{Related Work}
\subsection{UAV Absolute Visual Localization}
Recent UAV absolute localization research is motivated by the limitations of
navigation systems that depend on external signals or local motion integration:
GNSS may be unreliable or unavailable, while local odometry and SLAM
accumulate drift and do not directly recover global latitude-longitude
coordinates. A common alternative is image-based absolute localization, where
a UAV image is matched against a georeferenced satellite, orthomosaic, vector,
or 3D reference map \cite{liu2024review}. In contrast to relative visual
odometry, these methods aim to produce a global position estimate by anchoring
the observation to an external map.

Early systems performed this matching with classic, handcrafted local
features. Conte and Doherty registered onboard UAV video directly against a
georeferenced aerial image using local feature correspondences and
homography estimation, fusing the resulting position with inertial and
visual-odometry sensors to bound long-term drift \cite{conte2008integrated}.
Such approaches require no training data and integrate naturally with
classical filtering, but the underlying descriptors are fixed and
photometric: they degrade under seasonal, illumination, and viewpoint
differences between the UAV and reference imagery, and they struggle in
low-texture regions such as fields or water where few reliable keypoints
exist.

Convolutional neural networks (CNNs) later replaced handcrafted descriptors
with learned representations trained to be more invariant to the UAV-satellite
domain gap. Goforth and Lucey localized a UAV against pre-existing satellite
imagery by aligning CNN feature maps between the two views, jointly
optimizing consistency across adjacent UAV frames and the satellite
reference to improve accuracy \cite{goforth2019gpsdenied}. This moved the
matching signal from fixed gradient-based descriptors to representations
learned from data, but still framed localization as direct image-to-map
alignment rather than large-scale embedding retrieval, limiting scalability
to maps far larger than the immediate registration window.

More recently, deep-network retrieval has replaced direct image alignment as
the dominant paradigm, reframing localization as embedding-space search over
a large reference gallery rather than pairwise registration. Many such
learning-based UAV absolute localization methods have adopted a
retrieval-and-refinement paradigm. GLVL introduced a
global-local design in which a large-scale retrieval module first identifies a
similar map region and a fine-grained matching module then estimates a more
precise UAV coordinate \cite{li2023glvl}. This formulation is related to the
coarse-to-fine structure considered in this paper and is representative of
recent work that combines large-scale region retrieval with fine-grained local
matching.

DenseUAV focused on low-altitude urban self-positioning and emphasized the
importance of dense satellite sampling and localization-aware evaluation for
UAV-to-satellite matching \cite{dai2024denseuav}. Its benchmark and baseline
are useful for evaluating drone-satellite retrieval under urban appearance
changes, altitude variation, and dense candidate galleries. DenseUAV also
reports both retrieval and distance-sensitive evaluation, making it an
important reference point for UAV localization accuracy.

Recent alternatives also explore visual-inertial absolute positioning
\cite{tong2025rtapm}, vector-map localization \cite{hao2025vecmaploc}, and
multi-view visual indexing from rendered references \cite{li2025viewpoint}.
These approaches broaden the absolute-localization setting beyond pure
satellite-image retrieval, although they typically require additional sensor
streams, map modalities, or reference rendering pipelines.

\subsection{Vision Transformers and Cross-View Retrieval Architectures}
Vision Transformer (ViT) models replaced fixed convolutional receptive fields
with self-attention over image patches, allowing descriptors to aggregate
long-range spatial context directly from token interactions
\cite{dosovitskiy2021vit}, and early transformer retrieval studies showed
that ViT backbones optimized with metric-learning objectives produce compact,
discriminative global descriptors \cite{el2021vitretrieval}.
Self-supervised pretraining strengthened this direction: DINO and DINOv2
learn features with strong semantic organization that transfer across image
distributions without task-specific labels
\cite{caron2021dino,oquab2023dinov2}, and AnyLoc showed such
foundation-model features generalize across viewpoints and environments in
visual place recognition \cite{keetha2023anyloc}. These results motivate
using a strong pretrained backbone as a stable feature extractor and
adapting it to the target localization objective, as done in this work.

Applied directly to cross-view geo-localization, EgoTR modeled global
dependencies between street-view and aerial imagery with an evolving
transformer \cite{zhu2021egotr}, and TransGeo argued that a
transformer-only architecture can replace handcrafted cross-view alignment
operations by letting attention emphasize discriminative regions
\cite{zhu2022transgeo}. Attention thus allows image regions to be compared
through global context rather than only local convolutional patterns.

The strongest methods in our experimental comparison
(Tables~\ref{tab:soa_sues200}--\ref{tab:soa_denseuav}) build on these
backbones with explicit part- or region-level supervision. LPN partitions
the feature map into square-ring parts around the image center and
supervises each part separately, showing that surrounding context — not
only the central target — carries localization signal \cite{wang2021lpn}.
FSRA carries this idea to transformers: it ranks ViT tokens by activation
strength, segments them into semantic region groups, and aligns
corresponding regions across views, improving robustness to position and
scale shifts \cite{dai2022fsra}. CAMP combines a position-aware partition
branch, which captures fine-grained part features together with their
spatial arrangement, with a contrastive attributes mining strategy that
exploits attribute-level positives and negatives \cite{wu2024camp}. CEUSP,
the strongest published DenseUAV method, augments a compact ConvNeXt
backbone with a Rubik's-Cube attention module that rotates feature maps to
integrate multi-directional context, context-aware channel integration, and
a dynamic negative-sampling strategy \cite{xu2025ceusp}. A parallel
contrastive-sampling line — Sample4Geo, MCCG, DAC, and MEAN — is discussed
with the metric-learning literature below. Common to all these methods is
that discrimination among look-alike candidates is improved by supervising
\emph{where} content sits, not only what it is; the proposed pipeline
pursues the same goal through cluster-structured hard-negative mining on a
single pooled descriptor, plus an optional vector-quantized layout
re-ranking stage, rather than through per-part heads or alignment losses.

\subsection{Drone-Satellite Localization Challenges}
Drone-satellite localization is harder than standard image retrieval because
the query and reference images differ in viewpoint, altitude, scale, sensor
characteristics, illumination, and seasonal appearance. A UAV image may contain
oblique structures, motion blur, partial fields of view, or heading changes,
whereas the satellite reference is usually a top-down georeferenced raster.
These differences create a substantial domain gap, so a descriptor must remain
stable under geometric and photometric changes while still preserving enough
local detail to distinguish nearby map regions.

Large-map search further complicates the problem. Operational localization
does not usually provide a small closed-set gallery; instead, the system must
search many overlapping satellite crops extracted from a continuous map. This
creates severe class imbalance and many visually repetitive candidates, such as
roads, rooftops, fields, rivers, and industrial areas. The correct UAV footprint
may also overlap multiple reference tiles, making one-to-one retrieval labels an
imperfect training and evaluation target \cite{xu2024uavvisloc,li2025onetomany}.
These challenges motivate coarse-to-fine or clustered search designs that first
reduce the candidate set and then apply a more discriminative local retrieval
stage.

\subsection{Drone-Satellite Retrieval Benchmarks}
Drone-satellite retrieval benchmarks have shaped most current UAV
self-localization methods. SUES-200 introduced multi-height drone-satellite
matching with UAV images collected at four flight heights and corresponding
satellite views, highlighting altitude variation as a central challenge
\cite{zhu2023sues}. DenseUAV further moved the task toward low-altitude
self-positioning in urban environments \cite{dai2024denseuav}. ComplexUAV
extends the benchmark setting to high-altitude, dense, multi-terrain sampling
and evaluates contrastive learning, backbone choices, SAFA attention, and
generalization across multiple datasets \cite{chen2025complexuav}. These
benchmarks are usually formulated as retrieving a correct satellite tile or
class label from a prepared gallery.

Game4Loc introduced GTA-UAV, a game-engine UAV geo-localization benchmark
covering a continuous $81.3\,\mathrm{km}^2$ area with 33,763 drone images and
14,640 satellite tiles sampled at four zoom levels \cite{ji2024game4loc}. It
is particularly relevant to this work because it removes the strict
one-to-one pairing assumption and evaluates arbitrary drone--satellite pairs
through partial overlap, distance-sensitive metrics, and same-area/cross-area
training splits. Its satellite gallery can be viewed as a hierarchical tiling
reference: the same broad area is represented at several map zoom levels, so
retrieval must handle both scale changes and partial map overlap.

\begin{table*}[t]
\caption{Comparison of drone-satellite localization benchmarks.}
\label{tab:drone_satellite_benchmarks}
\centering
\scriptsize
\setlength{\tabcolsep}{4pt}
\begin{tabular}{@{}>{\raggedright\arraybackslash}p{0.15\textwidth}
                >{\raggedright\arraybackslash}p{0.25\textwidth}
                >{\raggedright\arraybackslash}p{0.30\textwidth}
                >{\raggedright\arraybackslash}p{0.22\textwidth}@{}}
\toprule
Dataset & Reference setting & Main variation & Task form \\
\midrule
University-1652 \cite{zheng2020university1652} &
Paired drone, satellite, and street-view images per building &
Multi-view supervision (drone/satellite/street) with synthetic multi-view drone rendering &
Closed-gallery multi-view retrieval \\
SUES-200 \cite{zhu2023sues} &
Paired drone and satellite views &
Multiple UAV flight heights &
Closed-gallery retrieval \\
DenseUAV \cite{dai2024denseuav} &
Dense urban satellite sampling &
Low-altitude urban viewpoint and appearance changes &
Dense tile retrieval and localization-aware evaluation \\
ComplexUAV \cite{chen2025complexuav} &
High-altitude, dense, multi-terrain sampling &
Terrain diversity, backbone choice, and cross-dataset generalization &
Benchmark retrieval across multiple datasets \\
UAV-VisLoc \cite{xu2024uavvisloc} &
Large satellite maps with UAV queries &
Diverse locations, platforms, altitudes, headings, and terrains &
Coordinate localization on georeferenced maps \\
Trim-UAV-\newline VisLoc \cite{li2025onetomany} &
Trimmed map tiles derived from UAV-VisLoc &
One-to-many imperfect overlap between query and references &
Overlap-aware retrieval with IoU information \\
GTA-UAV / \newline Game4Loc \cite{ji2024game4loc} &
Continuous synthetic map with four satellite zoom levels &
Altitude, attitude, scene, scale, and partial-overlap variation &
Weighted partial-match retrieval and localization in meters \\
\bottomrule
\end{tabular}
\end{table*}

This work mainly follows UAV-VisLoc and Trim-UAV-VisLoc because their problem
setting is closest to the intended operating condition of the proposed system.
UAV-VisLoc formulates UAV visual localization as finding real UAV coordinates
on a large satellite map rather than retrieving only a pre-cropped one-to-one
reference tile \cite{xu2024uavvisloc}. It provides 6,742 drone images and 11
satellite maps from diverse locations, terrains, UAV platforms, altitudes, and
headings. This makes it suitable for evaluating a pipeline that tiles a
GeoTIFF map, caches map descriptors, searches nearest map regions, and reports
geometric localization error rather than only closed-gallery classification
accuracy.

Trim-UAV-VisLoc is also central to this work because it exposes a practical
property of large-map localization: one UAV query may overlap several valid
reference crops, and the best tile is not always a single manually defined
class label. The dataset constructs a one-to-many imperfect matching setting
from UAV-VisLoc, adds IoU information, and introduces an Incomp-NPair loss to
better model overlapping references \cite{li2025onetomany}. This motivates the
adopted evaluation design: the satellite map is tiled into overlapping crops,
candidate retrieval is restricted by learned clusters, and a prediction is
judged by the Game4Loc IoU criterion (IoU $\geq 0.39$) instead of by a
strict one-to-one class label.

\subsection{Game4Loc Reference Protocol}
Game4Loc is used in this paper as a reference protocol for hierarchical
tiling, partial-match training, and area-split evaluation
\cite{ji2024game4loc}. Its GTA-UAV gallery contains satellite tiles from four
map zoom levels, which creates a multi-scale reference hierarchy. Rather than
assuming that each drone query corresponds to exactly one satellite tile, the
metadata lists positive and semi-positive satellite candidates for each drone
image. The official training condition uses metadata files such as
\texttt{same-area-drone2sate-train.json} or
\texttt{cross-area-drone2sate-train.json}; the default training mode
\texttt{pos\_semipos} samples from \texttt{pair\_pos\_sate\_img\_list} and
\texttt{pair\_pos\_semipos\_sate\_img\_list} with corresponding
overlap-derived positive weights.

The Game4Loc objective is a bidirectional weighted InfoNCE. Let
$\mathbf{u}_i$ and $\mathbf{v}_i$ be normalized drone and satellite
descriptors and let $\alpha_i$ denote the partial-match weight. Its
implementation converts this weight into a smoothing factor
$\epsilon_i=1-(1+\exp(-k\alpha_i))^{-1}$ and optimizes a mixture of a hard
diagonal contrastive target and a softened target:
\begin{equation}
\begin{split}
\mathcal{L}_{G4L} =
\frac{1}{2}\big[
&\mathcal{L}_w(\mathbf{U}\mathbf{V}^{T})+
\mathcal{L}_w(\mathbf{V}\mathbf{U}^{T})\big],\\
\mathcal{L}_w(\mathbf{S}) =
\frac{1}{B}\sum_i
&\left[
(1-\epsilon_i)\left(-S_{ii}+\log\sum_j e^{S_{ij}}\right)\right.\\
&\left.+\epsilon_i\left(-\frac{1}{B}\sum_j S_{ij}+
\log\sum_j e^{S_{ij}}\right)
\right].
\end{split}
\end{equation}
The same-area condition trains and tests on held-out samples from the same
geographic area, while the cross-area condition trains and tests on different
map regions. Evaluation is performed primarily in drone-to-satellite (D2S)
mode and reports Recall@K, average precision, spatial distance metric
(SDM@K), and mean distance error Dis@K. In the official evaluator, same-area
distance-threshold curves use 4 m increments up to 200 m, while cross-area
curves use 10 m increments up to 500 m.

\subsection{Metric Learning and Hard Negatives for Cross-View Localization}
Recent cross-view geo-localization methods increasingly rely on contrastive
metric learning rather than classification-only training. Sample4Geo demonstrated
that symmetric InfoNCE with hard negative sampling provides a strong
alternative to more complex aggregation and preprocessing pipelines
\cite{deuser2023sample4geo}. It shows that carefully designed sampling and
contrastive optimization can provide strong cross-view retrieval performance
without relying on complex preprocessing or aggregation modules.

ComplexUAV also reports strong performance from an InfoNCE-based contrastive
framework with dynamic distance sampling \cite{chen2025complexuav}. The
one-to-many UAV-VisLoc follow-up modifies the N-pair objective to handle
imperfect overlapping matches \cite{li2025onetomany}. These studies indicate
that hard-negative mining, multi-positive supervision, and overlap-aware
objectives are important for cross-view UAV localization.

A separate line of work strengthens the Sample4Geo/MCCG contrastive baseline
with additional spatial-feature supervision rather than sampling changes alone.
MCCG splits the backbone's spatial feature map into heatmap-ranked blocks and
applies part-level contrastive matching in addition to the global descriptor
objective \cite{shen2024mccg}. DAC builds directly on this codebase and adds a
domain/scene-alignment loss that matches flattened drone and satellite feature
maps across views, together with a distributionally-robust hard-negative
reweighting objective that penalizes difficult negative pairs more heavily
than plain InfoNCE \cite{xia2024dac}. MEAN extends the same design with a
combined cosine-similarity-and-MSE multilevel alignment loss for consistency
and invariance learning across viewpoints, while reusing DAC's hard-negative
reweighting objective, and reports the strongest published SUES-200 and
University-1652 numbers to date at a compact 36.5\,M-parameter ConvNeXt-Tiny
size \cite{chen2025mean}. Both methods indicate that supervising alignment at
the spatial-feature level, not only on the pooled global descriptor, is an
effective way to improve cross-view retrieval accuracy. The present work
instead supervises spatial structure only indirectly, through
cluster-consistency regularization on the pooled descriptor and structured
hard-negative mining, without an explicit spatial-feature alignment loss;
Section~\ref{sec:limitations} and the per-method comparison tables discuss
this trade-off further.

\subsection{Coarse-to-Fine Retrieval and Clustered Search}
Large-map localization is often too expensive to solve by exhaustive
fine-grained matching over every possible crop. Coarse-to-fine retrieval
therefore remains a common design pattern: a global representation narrows the
candidate set, and a more discriminative stage performs local ranking or
geometric refinement. Global-local UAV localization systems instantiate this
pattern with explicit refinement modules \cite{li2023glvl}, while dense-tile
benchmarks increase spatial recall by sampling many overlapping satellite
crops \cite{dai2024denseuav}. Approximate nearest-neighbor systems and FAISS
K-means provide a complementary route by organizing descriptors into searchable
partitions before fine ranking \cite{johnson2019billion}. These clustered
search strategies are particularly relevant when the reference map contains
many overlapping candidate tiles.

\section{Baseline Methods}
The proposed pipeline builds on three baseline method families commonly used
in visual
localization and image retrieval: Swin Transformer visual features,
FAISS K-means clustering, and
metric-learning losses. These baselines provide the feature extractor,
routing mechanism, and descriptor-learning objective used by the
proposed pipeline.

\subsection{Swin Transformer Backbone}
Swin Transformer introduces shifted-window self-attention, enabling local
attention computation while still exchanging information across windows
\cite{liu2021swin}. This property is well suited to map and UAV imagery, in which a
descriptor must preserve local texture cues such as roads, roofs, fields, and
shorelines while also encoding larger spatial layout. In the implemented
system, Swin-B is the primary backbone and ConvNeXt-T is supported as a
smaller alternative. The backbone is ImageNet-pretrained; its last stages
are unfrozen for fine-tuning, which is the only trainable retrieval
component since the raw normalized backbone output serves directly as the
descriptor. Keeping the earlier stages fixed reduces GPU memory use and
limits overfitting to satellite crop pairs.

\subsection{FAISS K-Means Routing}
K-means is employed as a routing mechanism rather than as a flat retrieval
method. The normalized backbone output is a unit descriptor, and FAISS provides an
efficient implementation of K-means and nearest-centroid search for large
embedding sets \cite{johnson2019billion}. This design is appropriate for large
GeoTIFF maps, as the map can be partitioned into candidate groups before
candidate scoring. The cluster assignments also create structured
hard-negative pools: samples from the same cluster are visually close and
thus informative for descriptor training. Persisting centroids further
stabilizes checkpoint resumption and cluster-consistency losses.

\subsection{Metric Learning}
The localization objective is naturally metric-based: matching crops are
expected to have high cosine similarity, whereas non-matching crops should
have lower similarity. The training objective therefore follows metric-learning
practice. Paired in-batch InfoNCE aligns anchor-positive descriptors, a
variance term prevents collapse, a same-cluster hard-negative margin suppresses
confusing alternatives, and a soft cluster-consistency loss encourages paired
crops to share K-means assignments. This follows the same principle as
contrastive and triplet metric learning
\cite{hadsell2006dimensionality,schroff2015facenet} while adapting the loss to
clustered routing over a large map.

\section{Problem Formulation}
Let $\mathcal{Q}=\{q_i\}$ be a set of UAV query images and
$\mathcal{M}=\{m_j\}_{j=1}^{N}$ be a set of map tiles extracted from a
georeferenced satellite image. For each query $q_i$, the objective is to
retrieve a ranked list of map tiles such that at least one tile overlaps the
ground-truth UAV location. The evaluation uses latitude and longitude
metadata from a sequence CSV file and transforms these coordinates to the map
coordinate reference system when required.

Instead of matching every query against all map tiles in a single flat search,
the proposed system first partitions map tiles into clusters in a learned
embedding space. A query searches one or more nearest clusters, and final
ranking is performed only over the candidate tiles in those clusters.

\section{Training Data}
\label{sec:training_data}
All models are trained with a cross-view drone-to-satellite (D2S)
supervision signal: each training sample pairs a UAV image with a satellite
crop of the same location, optionally carrying a continuous positive weight
consumed by the weighted InfoNCE objective. Every benchmark result reported
in this paper comes from a model trained, or fine-tuned, on that benchmark's
own official training split only — test locations are never seen during
training. Five training sources are used.

\paragraph{DenseUAV.}
The DenseUAV training split \cite{dai2024denseuav} contains 2,256 locations,
each with drone images captured at 80, 90, and 100\,m altitude and one
satellite tile per altitude. Training uses cross-altitude multi-positive
pairing: every drone image is paired with the same location's satellite tile
at \emph{all three} altitudes — all valid views of the same place — giving
nine pairs per location (20,304 pairs in total). Cross-altitude pairs are
softened by an altitude-gap weight $w=\exp(-|\Delta \mathrm{alt}|/\tau)$
with $\tau=30$\,m, so same-altitude pairs remain hard positives ($w{=}1$)
while $10$ and $20$\,m gaps receive $w{\approx}0.72$ and $0.51$.
Per-location GPS metadata is additionally loaded to exclude geographically
adjacent locations from the hard-negative pool.

\paragraph{SUES-200.}
The official 120-location training split of SUES-200 \cite{zhu2023sues}
provides drone images at four flight heights (150--300\,m), 50 images per
height per location, and one satellite image per location — 24,000 pairs,
all full-strength positives. The 80 test locations are held out entirely.

\paragraph{University-1652.}
The University-1652 training split \cite{zheng2020university1652} covers 701
buildings, each with $\sim$54 synthetic drone views and one satellite image.
Ten drone views are sampled per building per run, giving 7,010 pairs; all
drone views of a building share its satellite image as a full-strength
positive.

\paragraph{GTA-UAV (Game4Loc).}
The GTA-UAV same-area training metadata \cite{ji2024game4loc} is loaded in
the official \texttt{pos\_semipos} mode: each drone image is paired with
every positive (IoU $\geq 0.39$) and semi-positive (IoU $\geq 0.14$)
satellite tile listed for it, with each pair's overlap-derived value used
as its positive weight. The $\sim$27,000 same-area training drone images
expand to 123,095 weighted pairs under this scheme — the largest training
source used.

\paragraph{UAV-VisLoc with Game4Loc-style pairing.}
For the UAV-VisLoc experiments, the dataset \cite{xu2024uavvisloc} is
processed with the same Game4Loc tiling protocol \cite{ji2024game4loc}: a
pre-tiled satellite pyramid at four zoom levels serves as the gallery, and
for each drone image positive tiles are identified by geographic overlap
against the drone footprint (strong positive at IoU $\geq 0.39$,
semi-positive at IoU $\geq 0.14$, entering training with the IoU value as
its weight). The footprint is computed from altitude and heading metadata in
the per-sequence CSV files using a camera model with
$\text{FOV}_H = 36^\circ$ and $\text{FOV}_V = 52^\circ$. This source
provides real D2S signal over diverse terrain with authentic sensor
appearance differences.

\paragraph{Training transform.}
All training images are resized to the backbone input size ($384\times384$
for Swin-B, ConvNeXt-T, and ConvNeXt-B).
Augmentation includes color jitter, Gaussian blur or sharpening,
random cardinal rotation, JPEG compression simulation, and coarse dropout.
Random resized cropping is omitted to avoid removing overlapping content
that would corrupt the paired supervision signal.
Cluster-stat extraction at the start of each epoch uses a deterministic
transform: resize only, no augmentation.

\section{Method}
The proposed methodology is designed around the observation that UAV
localization over a large GeoTIFF map requires both scalable search and
fine-grained discrimination. A single global descriptor is efficient, but it
must separate many visually repetitive map regions in one embedding space. A
fully specialized local descriptor can improve discrimination, but applying it
over every map tile is computationally inefficient and difficult to cache. The
proposed architecture therefore uses a single retrieval descriptor — the
backbone's raw L2-normalized output embedding, adapted by fine-tuning the
last backbone stages — as both
the routing descriptor and the candidate-scoring descriptor. Training shapes
this embedding space to group visually and geographically
related crops, while evaluation uses nearest-cluster routing to reduce the
candidate set before scoring. This design adds no head parameters and keeps
all map embeddings cacheable.

\subsection{Backbone}
The visual encoder uses a pretrained backbone from \texttt{timm}
with the classification layer removed. The primary configuration is Swin-B at
$384\times384$ \cite{liu2021swin}; a ConvNeXt-T variant at the same input
resolution is supported for a $3\times$ smaller model with the identical
methodology. The backbone outputs
\begin{equation}
    \mathbf{h}=f_b(\mathbf{x})\in\mathbb{R}^{D},
\end{equation}
where $D=1024$ for Swin-B and $D=768$ for ConvNeXt-T. The backbone is
initialized with ImageNet-pretrained weights; the last backbone stages are
unfrozen for fine-tuning, which is the only mechanism by which the descriptor
is adapted.

\subsection{Retrieval Descriptor}
The retrieval descriptor is the backbone's own globally pooled output,
L2-normalized:
\begin{equation}
    \mathbf{c}(\mathbf{x}) =
    \mathrm{norm}\left(\mathbf{h}\right)\in\mathbb{R}^{D}.
\end{equation}
No projection head is used: the raw normalized backbone embedding serves
descriptor matching, K-means cluster assignment, cluster-consistency
training, and hard-negative mining alike, matching the head-free design used
by strong contrastive baselines \cite{deuser2023sample4geo,ji2024game4loc}.
Keeping the descriptor identical to the backbone output avoids an extra
projection whose benefit did not materialize in practice, and makes cached
gallery embeddings directly comparable across training stages. The number of
clusters is denoted by $K$ and is configurable in the implementation; the
default experimental setting uses $K=8$--$16$.

\subsection{Cluster-Structured Hard-Negative Mining}
\label{sec:hard_mining}
The K-means partition maintained for routing doubles as the training-time
negative-sampling structure. For each anchor, hard negatives are drawn
primarily from the anchor's own cluster — samples the current descriptor
already considers similar — and secondarily from the nearest neighboring
clusters, ranked by cosine similarity in the cached descriptor bank. A
remaining fraction is sampled globally at random to preserve coverage.

Two safeguards keep this mining signal clean. First, \emph{co-positive
exclusion}: any candidate that shares a positive satellite tile with the
anchor (including candidates linked through overlapping positives) is removed
from the negative pool, since pushing apart two views of the same place would
directly oppose the alignment objective. Second, the mining hardness is
scheduled: the negative loss is warmed up over the first training steps so
that early, unstable descriptors are not dominated by noisy hard negatives.
Because cluster assignments are refreshed from the current descriptor at each
rebuild, the mining pool tracks the model: as look-alike regions get pulled
apart, the clusters — and therefore the negatives — re-focus on whatever
confusions remain. Figure~\ref{fig:mining_comparison} contrasts this
cluster-structured sampling with the standard in-batch InfoNCE negative set.

\begin{figure*}[htbp]
\centering
\tikzset{
  negopen/.style = {circle, draw=red!70!black, thick, inner sep=1.6pt},
  negfill/.style = {circle, fill=red!70!black, inner sep=1.6pt},
  posdot/.style  = {circle, fill=green!55!black, inner sep=1.8pt},
  anchdot/.style = {star, star points=5, star point ratio=2.4, fill=black,
                    inner sep=1.5pt},
  graydot/.style = {circle, fill=black!30, inner sep=1.4pt},
  cent/.style    = {cross out, draw=blue!60!black, thick, inner sep=1.8pt},
  panel/.style   = {rounded corners=4pt, draw=black!60, thick},
  note/.style    = {font=\scriptsize, align=center},
}
\resizebox{0.92\textwidth}{!}{%
\begin{tikzpicture}

%% ── Panel (a): standard in-batch InfoNCE ────────────────────────────────
\begin{scope}
\draw[panel] (0,0) rectangle (7,5);
\node[font=\footnotesize\bfseries, anchor=south west] at (0.05,5.05)
  {(a) Standard in-batch InfoNCE negatives};

% anchor + positive
\node[anchdot] (a1) at (2.2,2.6) {};
\node[posdot]  at (2.55,2.95) {};

% hard region around the anchor: almost never sampled
\draw[dashed, black!50] (2.2,2.6) circle (0.85);
\node[note, black!60] at (2.2,1.35) {look-alike neighborhood:\\rarely sampled};

% random batch negatives, spread over the whole space (easy)
\foreach \p in {(5.9,4.3),(6.2,1.2),(4.8,0.7),(5.4,3.1),(6.5,2.5),
                (4.3,4.5),(3.9,1.0),(1.0,4.4),(0.8,0.9)}
  \node[negopen] at \p {};

\node[note] at (4.9,4.85) {negatives $=$ other batch items,\\drawn
  uniformly over the dataset};
\node[note, red!60!black] at (5.6,0.35)
  {distant $\Rightarrow$ easy $\Rightarrow$ vanishing gradient};
\end{scope}

%% ── Panel (b): cluster-structured mining ────────────────────────────────
\begin{scope}[shift={(8.6,0)}]
\draw[panel] (0,0) rectangle (7,5);
\node[font=\footnotesize\bfseries, anchor=south west] at (0.05,5.05)
  {(b) Cluster-structured mining (ours)};

% K-means cell boundaries (dashed) and centroids
\draw[dashed, black!45] (3.45,0) .. controls (3.2,1.8) and (3.6,3.4) .. (3.3,5);
\draw[dashed, black!45] (3.4,2.15) .. controls (4.8,2.3) and (5.8,2.0) .. (7,2.3);
\node[cent] at (1.9,2.7) {};
\node[cent] at (5.2,3.7) {};
\node[cent] at (5.3,1.0) {};
\node[note, black!55] at (0.85,4.72) {anchor's\\cluster};
\node[note, black!55] at (6.15,4.72) {nearest\\cluster};
\node[note, black!55] at (6.3,1.9) {other\\clusters};

% anchor + positive
\node[anchdot] at (2.0,2.4) {};
\node[posdot]  at (2.35,2.75) {};

% 50% same-cluster hard negatives, ranked by similarity (close to anchor)
\foreach \p in {(1.45,2.0),(2.65,1.85),(1.6,3.25),(2.85,2.95)}
  \node[negfill] at \p {};
\node[note, red!60!black] at (1.75,1.15) {50\,\% same cluster\\(hardest by
  similarity)};

% 30% nearest-cluster negatives
\foreach \p in {(4.35,3.4),(5.0,4.25),(4.7,2.75)}
  \node[negfill] at \p {};
\node[note, red!60!black] at (4.75,4.75) {30\,\% nearest\\cluster};

% 20% random global
\node[negfill] at (6.3,0.65) {};
\node[note, red!60!black] at (5.45,0.4) {20\,\% random};

% co-positive exclusion
\node[graydot] (cp) at (2.5,2.15) {};
\draw[red!70!black, thick] (2.36,2.01) -- (2.64,2.29);
\draw[red!70!black, thick] (2.36,2.29) -- (2.64,2.01);
\node[note, black!60, anchor=west] at (3.0,1.45)
  {co-positive: shares the anchor's\\positive tile $\to$ excluded};
\draw[-Stealth, black!50] (3.05,1.6) -- (2.62,2.05);
\end{scope}

%% ── Shared legend ────────────────────────────────────────────────────────
\begin{scope}[shift={(1.2,-0.75)}]
\node[anchdot]  at (0,0) {};   \node[note, anchor=west] at (0.15,0) {anchor};
\node[posdot]   at (2.0,0) {}; \node[note, anchor=west] at (2.15,0) {positive};
\node[negopen]  at (4.2,0) {}; \node[note, anchor=west] at (4.35,0) {in-batch negative};
\node[negfill]  at (7.3,0) {}; \node[note, anchor=west] at (7.45,0) {mined negative};
\node[cent]     at (10.2,0) {};\node[note, anchor=west] at (10.35,0) {K-means centroid};
\draw[dashed, black!45] (12.9,-0.05) -- (13.5,-0.05);
\node[note, anchor=west] at (13.6,0) {cluster boundary};
\end{scope}

\end{tikzpicture}%
}
\caption{Standard in-batch InfoNCE negatives versus the proposed
cluster-structured hard-negative mining, drawn in the retrieval embedding
space. \textbf{(a)}~With in-batch InfoNCE, the negative set is whatever else
happens to be in the batch — a uniform draw over the dataset. Most samples
land far from the anchor, where the contrastive gradient is already
vanishing, while the anchor's look-alike neighborhood (dashed) is almost
never sampled. \textbf{(b)}~Our mining uses the K-means partition maintained
for routing: 50\,\% of negatives come from the anchor's own cluster ranked by
descriptor similarity (the hardest confusers), 30\,\% from the nearest
neighboring cluster, and 20\,\% globally at random for coverage
(Section~\ref{sec:hard_mining}). Candidates sharing a positive tile with the
anchor (co-positives) are excluded so that two views of the same place are
never pushed apart. Because clusters are rebuilt from the current descriptor,
the sampling distribution re-focuses on whatever confusions remain as
training progresses.}
\label{fig:mining_comparison}
\end{figure*}

\subsection{Cluster Routing and Candidate Scoring}
Given L2-normalized K-means centroids
$\mathbf{C}=\{\mathbf{c}_1,\ldots,\mathbf{c}_{K}\}$, the routing score for an
image is computed as cosine similarity between its retrieval descriptor and all
centroids:
\begin{equation}
    \mathbf{s}(\mathbf{x}) =
    \mathbf{c}(\mathbf{x})\mathbf{C}^{T}\in\mathbb{R}^{K}.
\end{equation}
At evaluation time, these scores select the nearest map clusters. Candidate
tiles inside those clusters are scored with the same retrieval
descriptor and the selected scoring mode, such as cosine, centered cosine,
cluster-centered cosine, or cluster-Mahalanobis scoring.

\begin{figure*}[htbp]
\centering
\tikzset{
  block/.style = {draw, rounded corners=2pt, align=center, font=\small,
                  minimum width=3.6cm, minimum height=0.62cm},
  arr/.style   = {-Stealth, semithick},
  vqarr/.style = {-Stealth, semithick, teal!70!black},
}
% semantic-map cell colors (schematic: vegetation/road/building/water/roof)
\colorlet{cG}{green!55!black!55}
\colorlet{cR}{black!35}
\colorlet{cB}{orange!75}
\colorlet{cW}{cyan!60!blue!55}
\colorlet{cO}{red!65}
% \semmap{x}{y}{cell size}{rows}: 6x6 colored ID map, top-left at (x,y)
\def\semmap#1#2#3#4{%
  \begin{scope}[shift={(#1,#2)}]
    \foreach \row [count=\yi from 0] in {#4} {
      \foreach \c [count=\xi from 0] in \row {
        \path[fill=\c, draw=white, line width=0.25pt]
          ({\xi*#3}, {-\yi*#3-#3}) rectangle ++(#3,#3);
      }}
    \draw[black!60, line width=0.4pt] (0, {-6*#3}) rectangle ({6*#3}, 0);
  \end{scope}}
\resizebox{0.82\textwidth}{!}{%
\begin{tikzpicture}

%% ─── Backbone (left column) ──────────────────────────────────────────────

\node[block, fill=gray!12]                       (inp)  {Input $[B,3,384,384]$};
\node[block, fill=teal!10,   below=3mm of inp]   (emb)  {Patch Embed ($4{\times}4$, stride~4)};
\node[font=\scriptsize,gray, right=2mm of emb]   {$[B,96^2,C{=}128]$};

\node[block, fill=cyan!8,    below=6mm of emb]   (s0)   {Stage~0 \quad $\times$2 blocks};
\node[font=\scriptsize,gray, right=2mm of s0]    {$96{\times}96,\ C{=}128$};

\node[block, fill=cyan!14,   below=6mm of s0]    (s1)   {Stage~1 \quad $\times$2 blocks};
\node[font=\scriptsize,gray, right=2mm of s1]    {$48{\times}48,\ C{=}256$};

\node[block, fill=cyan!20,   below=6mm of s1]    (s2)   {Stage~2 \quad $\times$18 blocks};
\node[font=\scriptsize,gray, right=2mm of s2]    {$24{\times}24,\ C{=}512$};

\node[block, fill=cyan!26,   below=6mm of s2]    (s3)   {Stage~3 \quad $\times$2 blocks};
\node[font=\scriptsize,gray, right=2mm of s3]    {$12{\times}12,\ C{=}1024$};

\node[block, fill=teal!8,    below=6mm of s3]    (avgp) {Global AvgPool};
\node[font=\scriptsize,gray, right=2mm of avgp]  {$[B,1024]$};

\node[block, fill=violet!15, below=6mm of avgp]  (l2)   {L2-normalize};

\node[block, fill=gray!20,   below=6mm of l2,
      line width=1.2pt]                           (ret)  {\textbf{Retrieval descriptor}\\$[B,1024]$};

\draw[arr] (inp) -- (emb);
\draw[arr] (emb) -- (s0);
\draw[arr] (s0)  -- node[font=\tiny,gray,right]{patch merge} (s1);
\draw[arr] (s1)  -- node[font=\tiny,gray,right]{patch merge} (s2);
\draw[arr] (s2)  -- node[font=\tiny,gray,right]{patch merge} (s3);
\draw[arr] (s3)  -- (avgp);
\draw[arr] (avgp)-- (l2);
\draw[arr] (l2)  -- (ret);

%% ─── VQ semantic re-ranking panel (right) ────────────────────────────────

% quantization box, fed from the stage-3 spatial grid
\node[block, fill=teal!8, minimum width=4.6cm] (vqquant) at (8.6,-5.6)
  {nearest-centroid quantization\\
   \scriptsize shared codebook $\mathbf{B}\in\mathbb{R}^{K\times C}$
   ($k$-means over gallery cells)};
\draw[vqarr, dashed] (s3.east) to[out=0, in=200]
  node[below, sloped, font=\tiny, teal!60!black]
  {stage-3 grid $12{\times}12{\times}1024$} (vqquant.west);

% query + candidate semantic maps
\semmap{6.6}{-0.6}{0.30}{{cG,cG,cR,cB,cB,cG},{cG,cO,cR,cB,cB,cG},{cR,cR,cR,cR,cR,cR},{cB,cB,cR,cG,cG,cG},{cB,cB,cR,cG,cW,cW},{cG,cG,cR,cW,cW,cW}}
\node[font=\scriptsize, align=center] at (7.5,-0.25) {query map $Q$};
\semmap{9.6}{-0.6}{0.30}{{cB,cB,cB,cR,cG,cG},{cG,cB,cB,cR,cO,cG},{cR,cR,cR,cR,cR,cR},{cW,cG,cG,cR,cB,cB},{cW,cW,cG,cR,cB,cB},{cW,cW,cG,cR,cG,cG}}
\node[font=\scriptsize, align=center] at (10.5,-0.25) {candidate map $T$};
\node[font=\tiny, gray, align=center] at (9.0,-2.75)
  {$H{\times}W$ centroid IDs — 144\,B/image};
\draw[vqarr] (vqquant.north) to[out=90, in=270] (7.5,-2.5);
\draw[vqarr] (vqquant.north) to[out=90, in=270] (10.5,-2.5);

% rotation search row
\node[font=\scriptsize, align=left, anchor=west] at (6.4,-3.25)
  {rotation search over $\mathrm{rot}_r(T)$:};
\semmap{6.6}{-3.6}{0.15}{{cB,cB,cB,cR,cG,cG},{cG,cB,cB,cR,cO,cG},{cR,cR,cR,cR,cR,cR},{cW,cG,cG,cR,cB,cB},{cW,cW,cG,cR,cB,cB},{cW,cW,cG,cR,cG,cG}}
\semmap{7.9}{-3.6}{0.15}{{cG,cG,cR,cB,cB,cG},{cG,cO,cR,cB,cB,cG},{cR,cR,cR,cR,cR,cR},{cB,cB,cR,cG,cG,cG},{cB,cB,cR,cG,cW,cW},{cB,cG,cR,cW,cW,cW}}
\semmap{9.2}{-3.6}{0.15}{{cG,cG,cR,cG,cW,cW},{cB,cB,cR,cG,cW,cW},{cB,cB,cR,cG,cG,cW},{cR,cR,cR,cR,cR,cR},{cG,cO,cR,cB,cB,cG},{cG,cG,cR,cB,cB,cB}}
\semmap{10.5}{-3.6}{0.15}{{cW,cW,cW,cR,cG,cB},{cW,cW,cG,cR,cB,cB},{cG,cG,cG,cR,cB,cB},{cR,cR,cR,cR,cR,cR},{cG,cB,cB,cR,cO,cG},{cG,cB,cB,cR,cG,cG}}
\node[font=\tiny] at (7.05,-4.75) {$0^\circ$: 0.25};
\node[font=\tiny, teal!60!black] at (8.35,-4.75) {$90^\circ$: \textbf{0.97}};
\node[font=\tiny] at (9.65,-4.75) {$180^\circ$: 0.28};
\node[font=\tiny] at (10.95,-4.75) {$270^\circ$: 0.14};
\draw[teal!70!black, line width=0.9pt, rounded corners=1pt]
  (7.82,-4.58) rectangle (8.88,-3.52);

% blended score
\node[block, fill=gray!10, minimum width=4.6cm] (vqscore) at (8.6,-7.4)
  {$m(q,t)=\max_r \frac{1}{HW}\sum_p
     \mathbb{1}[Q{=}\mathrm{rot}_r(T)] = 0.97$\\[1pt]
   \scriptsize $S'(q,t)=S(q,t)+\alpha\,m(q,t)$ \quad (top-$N$ only)};
\draw[vqarr] (8.6,-4.9) -- (vqscore.north);

%% ─── Backbone bounding box ───────────────────────────────────────────────

\begin{scope}[on background layer]
  \node[draw=teal!55, dashed, rounded corners=5pt,
        inner sep=8pt, fill=teal!3,
        fit=(inp)(emb)(s0)(s1)(s2)(s3)(avgp),
        label={[font=\footnotesize\itshape, color=teal!65]above:
          {Swin-B Backbone (last stages fine-tuned)}}] (bbbox) {};
\end{scope}

\end{tikzpicture}%
}
\caption{Model architecture and vector-quantized semantic re-ranking.
\textbf{Left:} Swin-B backbone with four stages of progressively downsampled
spatial tokens; the globally pooled, L2-normalized 1024-D backbone output is
the retrieval descriptor — no projection head. The same descriptor serves
matching, K-means cluster routing, candidate scoring, and the
cluster-structured hard-negative mining of Section~\ref{sec:hard_mining}.
\textbf{Right} (Section~\ref{sec:vq_rerank}): the stage-3 spatial grid of
each image is quantized against one shared $k$-means codebook into an
$H{\times}W$ map of centroid IDs (schematic $6{\times}6$ shown; actual
$12{\times}12$, 144 bytes per image). A query's top-$N$ coarse candidates are
re-scored by mean cell-wise ID agreement $m(q,t)$, maximized over the four
cardinal rotations of the candidate map (here $90^\circ$ wins with
$m{=}0.97$; reflections deliberately excluded), and blended into the coarse
similarity within the top-$N$ only.}
\label{fig:architecture}
\end{figure*}

\subsection{Methodological Comparison with Close Baselines}
The proposed method is closest to GLVL, DenseUAV-style dense retrieval,
Sample4Geo, and TransGeo, but it occupies a different design point from each.
GLVL and the proposed method both adopt a coarse-to-fine localization
philosophy. GLVL retrieves a coarse map region and then applies a local
matching module for precise coordinate estimation \cite{li2023glvl}. The
proposed method instead keeps the refinement stage in descriptor space: the
retrieval descriptor selects candidate regions by learned embedding clusters and scores
map tiles inside those clusters. This avoids invoking a
separate local-feature matcher for every query and makes map-side computation
straightforward to cache.

DenseUAV-style retrieval increases localization resolution by constructing a
dense satellite gallery and evaluating UAV queries against many candidate
tiles \cite{dai2024denseuav}. This strategy is effective when the candidate
gallery is fixed and densely annotated, but its online cost grows with the
number of satellite crops. The proposed method uses the Game4Loc multi-zoom tile gallery,
but reduces the number of online comparisons by
searching only the top-$K_c$ embedding clusters. Thus, the main comparison to
DenseUAV is not only Recall@K, but also candidates per query and latency per
query.

Sample4Geo provides a strong contrastive-learning baseline because it learns a
single global retrieval descriptor with symmetric InfoNCE and hard-negative
sampling \cite{deuser2023sample4geo}. The proposed method inherits the same
general metric-learning motivation but changes the retrieval structure. Rather
than performing flat global ranking over all tiles, the proposed architecture
optimizes one retrieval descriptor for stable routing and uses the same
space for local candidate scoring. Clustered search reduces the number of
candidate comparisons while keeping map embeddings cacheable.

TransGeo is a relevant transformer-based cross-view retrieval baseline because
it uses attention to model long-range correspondences between views
\cite{zhu2022transgeo}. In contrast, the proposed method does not train a full
transformer retrieval model end to end. It fine-tunes only the last stages of
a pretrained backbone, with the raw normalized output as the descriptor. This
choice reduces training cost and makes the
comparison focus on the proposed cluster-routed retrieval and mining
mechanism rather than on backbone capacity alone.

\section{Training Objective}
\subsection{Descriptor Training}
Training fine-tunes the last backbone stages; the descriptor itself is the
normalized backbone output, so no separate head parameters exist. The
objective is to assign
anchor-positive pairs to the same K-means cluster while maintaining descriptor
variance and suppressing hard negatives.

For a mini-batch of retrieval descriptors $\mathbf{q}_i$ and
$\mathbf{p}_i$, paired in-batch InfoNCE is
\begin{equation}
\mathcal{L}_{align}
=-\frac{1}{B}\sum_i
\log
\frac{\exp(\mathbf{q}_i^T\mathbf{p}_i/\tau)}
{\sum_j\exp(\mathbf{q}_i^T\mathbf{p}_j/\tau)},
\quad \tau=0.10.
\end{equation}

The variance regularizer prevents collapse:
\begin{equation}
\mathcal{L}_{var}=
\frac{1}{D}\sum_d
\max\left(0,\frac{1}{\sqrt{D}}-\sqrt{\mathrm{Var}(\mathbf{Z}_{:,d})+10^{-4}}\right),
\end{equation}
where $\mathbf{Z}$ stacks anchor and positive descriptors and $D$ is the
descriptor dimension ($1024$ for Swin-B, $768$ for ConvNeXt-T).

For sampled negatives $\mathbf{n}_{ij}$, the same-cluster hard-negative margin
is
\begin{equation}
\mathcal{L}_{neg} =
\mathbb{E}_{i,j}\left[
\max\left(0,
\frac{1}{2}(\mathbf{q}_i^T\mathbf{n}_{ij}
+\mathbf{p}_i^T\mathbf{n}_{ij})-0.60
\right)^2
\right].
\end{equation}

The implemented training loss is
\begin{equation}
\mathcal{L} =
1.0\mathcal{L}_{align}
+1.0\mathcal{L}_{var}
+0.2\mathcal{L}_{neg}
+2.0\mathcal{L}_{consistency}.
\end{equation}
The cluster-consistency term is activated once centroids are available. It
matches the soft centroid assignment distributions of each anchor-positive
pair:
\begin{equation}
\mathcal{L}_{consistency} =
\frac{1}{2}\left[
\mathrm{KL}(\pi(\mathbf{q}_i)\|\pi(\mathbf{p}_i))+
\mathrm{KL}(\pi(\mathbf{p}_i)\|\pi(\mathbf{q}_i))
\right],
\end{equation}
where $\pi(\mathbf{x})=\mathrm{softmax}(\mathbf{x}\mathbf{C}^{T}/0.10)$.
The primary monitoring metric is the anchor-positive same-cluster probability.

\subsection{Negative Sampling}
Following the cluster-structured mining of Section~\ref{sec:hard_mining},
training samples negatives using a mixture of 50\% same-cluster, 30\% nearby
cluster, and 20\% random global negatives. Hard same-cluster samples are
selected by similarity in the cached descriptor bank, and candidates sharing
a positive tile with the anchor are excluded from the pool.

\subsection{Dead-Cluster Filtering}
During training, the trainer identifies unreliable routing clusters using
anchor-positive assignment consistency rather than cluster compactness. High
intra-cluster similarity is expected because the descriptor is explicitly
trained to group similar terrain and corresponding crop pairs. A cluster is
therefore treated as dead only when its anchor-positive same-cluster rate is
below 0.50, or when the cluster descriptors are nearly collapsed according to
very low per-dimension descriptor variance. These dead clusters are recorded
as a diagnostic property of the descriptor space and saved in the checkpoint
together with representative sample indices. The current training procedure
does not remove those samples; the diagnostic is used for inspection.

\section{Clustering and Checkpointing}
At the beginning of each epoch, the trainer extracts retrieval
descriptors for anchors and positives, runs FAISS K-means
\cite{macqueen1967methods,johnson2019billion},
assigns samples to clusters, and computes the anchor-positive same-cluster
probability. K-means is re-run when this probability is below 0.80. The first
epoch uses multiple random restarts, while later epochs warm-start from the
previous centroids. Checkpoints store the backbone,
optimizer state, loss history, training parameters, K-means
centroids, diagnostic per-cluster mean embeddings, and the last training mode.
Checkpoint filenames are keyed by backbone, for example
\texttt{latest\_swin\_b\_nohead.pt} or \texttt{latest\_convnext\_t.pt}, with
timestamped backups saved before overwriting.

\begin{figure}[htbp]
\centering
\begin{minipage}{0.95\linewidth}
\begin{algorithmic}[1]
\STATE Initialize pretrained backbone (Swin-B or ConvNeXt-T)
\STATE Freeze early backbone stages; unfreeze the last stages for fine-tuning
\FOR{each training epoch}
    \STATE Extract L2-normalized retrieval descriptors
    \STATE Rebuild or warm-start FAISS K-means when same-cluster probability is low
    \STATE Mine hard negatives from same and nearest clusters (co-positive exclusion)
    \STATE Train with InfoNCE, variance, negative, and consistency losses
\ENDFOR
\STATE Save model weights, K-means centroids, training parameters, and diagnostics
\end{algorithmic}
\end{minipage}
\caption{Descriptor training procedure.}
\label{alg:training}
\end{figure}

\section{Evaluation Pipeline}
\subsection{Map Tiling}
\label{sec:eval}
The evaluation gallery follows the Game4Loc tiling protocol
\cite{ji2024game4loc}.
Each area in UAV-VisLoc is accompanied by a pre-tiled satellite image pyramid:
the full-resolution map is downsampled and sliced into non-overlapping
$256\times256$-pixel tiles at multiple zoom levels.
Given the full-resolution satellite image of dimensions $H\times W$ pixels,
the maximum zoom level is
\begin{equation}
    z_{\max} = \left\lceil \log_2 \frac{\max(H,W)}{256} \right\rceil.
\end{equation}
At zoom level $z$, the image is scaled by factor $s = 2^{z_{\max}-z}$,
yielding a scaled canvas of size
$\lceil H/s \rceil \times \lceil W/s \rceil$ pixels.
This canvas is partitioned into a grid of $256\times 256$ tiles indexed by
column $x$ and row $y$.
Coarser zoom levels ($z < z_{\max}$, larger $s$) produce tiles that cover
proportionally larger geographic footprints, spanning from tens of metres at
fine zoom to several hundred metres at coarse zoom.
The gallery spans four consecutive zoom levels to present tiles at a range
of geographic scales.

Each tile is identified by the naming convention
\texttt{\{area\}\_\{z\}\_\{x\}\_\{y\}.png}.
The geographic centre and bounds of a tile are recovered analytically from
the known area bounding box $(lat_{lt}, lon_{lt}, lat_{rb}, lon_{rb})$:
\begin{equation}
\begin{split}
    lat_c &= lat_{lt} -
        \frac{(y+\tfrac{1}{2})\cdot 256}{\lceil H/s \rceil}
        (lat_{lt} - lat_{rb}),\\
    lon_c &= lon_{lt} +
        \frac{(x+\tfrac{1}{2})\cdot 256}{\lceil W/s \rceil}
        (lon_{rb} - lon_{lt}).
\end{split}
\end{equation}
No online GeoTIFF slicing is required; tile images ship with the dataset.

\paragraph{Positive tile labelling.}
For each drone query, positives tiles are identified by pixel-space IoU
between the tile polygon and the UAV ground footprint.
The footprint is computed from the recorded altitude, heading angle, and
the camera field of view ($\mathrm{FOV}_H=36^\circ$, $\mathrm{FOV}_V=52^\circ$)
following the Game4Loc geometric model \cite{ji2024game4loc}.
Both the footprint and tile corners are projected into the scaled pixel space
at zoom $z$ before the intersection is computed.
A tile is labelled a \emph{positive} if $\mathrm{IoU} \geq 0.39$, a
\emph{semi-positive} if $\mathrm{IoU} \geq 0.14$, and a \emph{negative}
otherwise. A tile that falls fully within the drone footprint
(tile recall $\geq 0.80$) is also counted as a positive even when the
IoU falls below the primary threshold; this prevents fine-zoom tiles
from being incorrectly excluded when the drone footprint is much larger
than the tile.

Tile embeddings are cached keyed by the checkpoint hash, so repeated
evaluations over the same area and checkpoint skip gallery extraction entirely.
Query embeddings are cached separately per sequence.

\subsection{Coarse-to-Fine Retrieval}
Map tiles are embedded with the retrieval descriptor and clustered with FAISS. The
cluster data cache stores assignments, centroids, and a cluster-to-tile index.

For a query image, the retrieval descriptor finds the top-$K_c$ nearest
clusters. The candidate set is the union of all tiles assigned to those
clusters. Candidate tiles are scored directly with the retrieval descriptor
and the configured scoring mode. The number of searched clusters, $K_c$,
controls the trade-off between candidate set size and recall.

The online scoring cost is reduced in proportion to the fraction of clusters
searched. Let $N$ be the number of gallery tiles, $K$ the number of routing
clusters, and $K_c$ the number of nearest clusters searched per query. If tile
assignments are approximately balanced, the expected candidate count is
\begin{equation}
    \mathbb{E}[|\mathcal{C}(q)|] \approx \frac{K_c}{K}N,
\end{equation}
so the fine-stage descriptor scoring cost is reduced by a factor of
$K/K_c$ compared with exhaustive gallery scoring. For example, with $K=16$
and $K_c=2$, the expected fine-stage candidate set contains $N/8$ tiles.

The recall loss introduced by routing depends on whether the positive tile is
assigned to one of the searched clusters. Let $p_s$ denote the probability
that a positive query-tile pair is recovered by one searched cluster. Under a
simple independent-cluster approximation, the coarse-stage miss probability
after searching $K_c$ clusters is
\begin{equation}
    P_{\mathrm{miss}} \approx (1-p_s)^{K_c}.
\end{equation}
Thus, if $p_s=0.90$ and $K_c=2$, then
$P_{\mathrm{miss}}\approx(1-0.90)^2=0.01$, meaning that the coarse stage
retains the positive candidate with approximately 99\% probability while
scoring only one eighth of the gallery.

\subsection{Vector-Quantized Semantic Re-Ranking}
\label{sec:vq_rerank}
On densely tiled galleries, coarse retrieval exhibits a characteristic gap:
Recall@10 approaches saturation while Recall@1 remains several points lower,
i.e., the correct tile is almost always among the strongest candidates but is
not always ranked first. This gap has a structural cause. When the tile
sampling interval is much smaller than the tile footprint — DenseUAV samples
$\sim$100\,m tiles at $\sim$20\,m intervals — a query's strongest candidates
are near-duplicate, heavily overlapping views of the same ground: they
contain nearly the same buildings, roads, and vegetation, differing mainly
in \emph{where} that content sits within the frame. The globally pooled
retrieval descriptor is structurally blind to exactly this difference:
average pooling discards the spatial arrangement of features, so two
shifted views of the same content map to nearly identical descriptors
regardless of how well the descriptor is trained. Separating such
candidates therefore requires a second, \emph{spatially structured}
representation that preserves arrangement. An optional final stage
re-orders each query's top-$N$ coarse candidates using vector-quantized
semantic maps derived from intermediate backbone features — retaining the
layout information that pooling destroys, at a per-image cost of a few
hundred bytes.

This mechanism also predicts where the stage applies: on galleries whose
tiles are distinct, non-overlapping locations (one tile per location, or
one satellite image per building), coarse errors occur between
different-content places that the global descriptor already separates, and
layout agreement contributes no additional signal — the re-ranking stage
should then be disabled. Table~\ref{tab:soa_denseuav} reflects this:
the re-ranked rows are reported for the densely tiled DenseUAV benchmark,
where the stage contributes its largest gains.

Let the stage-3 feature grid of an image be an $H \times W$ array of cell
descriptors $\mathbf{x}_{p} \in \mathbb{R}^{C}$ (for Swin-B at $384^2$,
$H{=}W{=}12$ and $C{=}1024$). A single codebook
$\mathbf{B} \in \mathbb{R}^{K \times C}$ is fit once by $k$-means over
$\ell_2$-normalized cell descriptors sampled across the gallery, following
the vector-quantization tradition of visual-word retrieval
\cite{sivic2003video} and product quantization \cite{jegou2011product}.
Every image is then stored as an $H \times W$ array of nearest-centroid
indices — a \emph{semantic map} occupying $H\!\cdot\!W$ bytes per image for
$K \le 256$ (144 bytes at stage 3, a $\sim$4000$\times$ reduction over the
full-precision grid), plus one shared codebook.

The agreement between a query map $Q$ and a candidate map $T$ is the mean
cell-wise identity, maximized over the four cardinal rotations of the
candidate:
\begin{equation}
    m(q,t) = \max_{r \in \{0^\circ,90^\circ,180^\circ,270^\circ\}}
    \frac{1}{HW} \sum_{p} \mathbb{1}\!\left[\,Q(p) = \mathrm{rot}_r(T)(p)\,\right].
\end{equation}
The rotation search accounts for arbitrary UAV yaw against north-up satellite
tiles; reflections are deliberately excluded so that mirrored look-alike
structures — a recurring false-match mode in dense urban galleries — score
low rather than high. Exact index identity is used rather than soft
centroid-similarity credit: spatially adjacent tiles share similar content
whose cells quantize to \emph{neighboring} clusters, and partial credit
rewards precisely those near-miss cells, whereas consistent exact-identity
agreement is delivered only by the true location.

The final score blends the map agreement into the coarse similarity within
each query's top-$N$ candidates only,
\begin{equation}
    S'(q,t) = S(q,t) + \alpha\, m(q,t), \qquad t \in \mathrm{top}\text{-}N(q),
\end{equation}
with all other gallery scores unchanged. Because the additive term is
non-negative and confined to the top-$N$ set, no candidate outside the set
can be displaced into it, so Recall@$N$ is invariant by construction and the
stage acts purely on the ordering within the shortlist. The added cost is
$O(N \cdot HW)$ integer comparisons per query — negligible against the
embedding forward pass — and the per-image payload is small enough to hold
the semantic maps of an entire city-scale gallery in a few megabytes.

\subsection{Comparison with Alternative Coarse-to-Fine Strategies}
The proposed retrieval mechanism differs from several common coarse-to-fine
localization strategies. In global-local retrieval systems such as GLVL, the
coarse stage retrieves a candidate map region and the fine stage applies a
local matching module to refine the coordinate estimate \cite{li2023glvl}.
This approach provides explicit spatial refinement, but it requires a
separate fine-localization module after retrieval. In contrast, the proposed
method restricts the candidate set in descriptor space and then scores those
candidates using the same descriptor.

Dense tiling methods, including DenseUAV-style benchmark protocols, increase
localization resolution by evaluating many overlapping satellite crops
\cite{dai2024denseuav}. This strategy is simple and map-geometry compatible,
but the search cost grows with tile density. The proposed method uses the Game4Loc pre-tiled gallery at multiple zoom
levels, yet reduces the candidate set through learned embedding clusters
before final scoring. Classical image-retrieval pipelines
often use a global descriptor followed by geometric verification with local
features; such methods can improve precision but add feature matching,
outlier rejection, and pose-estimation costs. By comparison, the proposed
pipeline keeps both stages descriptor-based and therefore directly cacheable
for large map tiles.

Recent absolute localization methods also incorporate temporal or
multi-sensor refinement, such as visual-inertial optimization in RTAPM
\cite{tong2025rtapm}, vector-map matching in VecMapLocNet
\cite{hao2025vecmaploc}, and rendered-view indexing from 3D references
\cite{li2025viewpoint}. These methods can improve robustness when additional
signals or map modalities are available. The proposed method is restricted to
image-to-GeoTIFF matching, making it applicable to settings where only UAV
imagery, a georeferenced raster map, and query metadata are available.

\begin{table}[htbp]
\caption{Comparison of Coarse-to-Fine Localization Strategies}
\label{tab:coarse_fine_comparison}
\centering
\footnotesize
\setlength{\tabcolsep}{3pt}
\begin{tabular}{>{\raggedright\arraybackslash}p{0.23\columnwidth}
                >{\raggedright\arraybackslash}p{0.32\columnwidth}
                >{\raggedright\arraybackslash}p{0.34\columnwidth}}
\toprule
Strategy & Coarse stage & Fine stage \\
\midrule
Global-local retrieval \cite{li2023glvl} &
Retrieve candidate region with a global descriptor &
Apply local matching or coordinate refinement \\
Dense tile retrieval \cite{dai2024denseuav} &
Enumerate many overlapping map crops &
Rank dense crop candidates by image similarity \\
Game4Loc hierarchical tiling \cite{ji2024game4loc} &
Construct a multi-zoom satellite gallery and train with partial-match weights &
Rank drone-to-satellite candidates and score with Recall, AP, SDM, and distance error \\
Temporal or multi-sensor absolute localization \cite{tong2025rtapm,hao2025vecmaploc,li2025viewpoint} &
Use retrieval, vector maps, rendered views, or sensor priors &
Refine with temporal filtering, map constraints, or optimization \\
Proposed cluster-routed retrieval &
Select top-$K_c$ learned embedding clusters &
Score candidates with the retrieval descriptor; optional VQ semantic re-ranking of the top candidates \\
\bottomrule
\end{tabular}
\end{table}

\subsection{Ground-Truth Filtering and Query Orientation}
For UAV-VisLoc sequence 01, query geometry is read from \texttt{01.csv}.
Only queries whose ground-truth coordinate falls inside the selected GeoTIFF
bounds are evaluated. If enabled, query orientation metadata can be used as a
preprocessing correction before model inference; otherwise the query is
embedded directly and scored against the selected map-tile candidates.

\subsection{Metrics}
Two hit conventions are used, one per benchmark family. On the open-map
benchmarks (UAV-VisLoc and GTA-UAV), evaluation follows the Game4Loc
IoU-based matching protocol \cite{ji2024game4loc}: a retrieved tile $m_j$
is counted as a hit for query $q_i$ if and only if it is a positive tile
under the criterion defined in Section~\ref{sec:eval}:
\begin{equation}
  \mathrm{hit}(q_i, m_j) =
  \begin{cases}
    1 & \text{if } \mathrm{IoU}(m_j, F_i) \geq 0.39
        \text{ or } \mathrm{recall}(m_j, F_i) \geq 0.80, \\
    0 & \text{otherwise,}
  \end{cases}
\end{equation}
where $F_i$ is the drone ground footprint computed from altitude, heading,
and camera FOV ($\mathrm{FOV}_H=36^\circ$, $\mathrm{FOV}_V=52^\circ$), and
$\mathrm{recall}(m_j, F_i) = |m_j \cap F_i| / |m_j|$ is the fraction of
the tile area covered by the drone footprint.
Recall@K is then
\begin{equation}
\mathrm{Recall@K} =
\frac{1}{|\mathcal{Q}|}\sum_i
\mathbb{1}\!\left[
\exists m_j\in \mathrm{TopK}(q_i):
\mathrm{hit}(q_i, m_j) = 1
\right].
\end{equation}
Average precision (AP) is computed per query from the binary hit labels
on the ranked retrieved list.
SDM@K weights the top-$K$ retrieved locations by both rank and
geographic closeness:
\begin{equation}
\mathrm{SDM@K} =
\frac{\sum_{r=1}^{K}(K-r+1)\exp(-\gamma d_r)}
{\sum_{r=1}^{K}(K-r+1)},
\end{equation}
where $d_r$ is the distance between the query location and the tile retrieved
at rank $r$, and $\gamma$ is the evaluator's distance scale. The evaluator
reports \texttt{sdm\_at\_1} and \texttt{sdm\_at\_k}. It also reports
\texttt{dis\_at\_1} and \texttt{dis\_at\_k}, the mean geographic error in
meters for the top-1 result and the top-$K$ retrieved set. Query throughput is
reported as FPS after cached map preprocessing.

On the closed-gallery benchmarks, the hit criterion is instead an exact
identity match: a retrieved tile of the query's location ID on DenseUAV, or
of the query's class on SUES-200 and University-1652. Recall@K and AP are
computed from these binary labels with the same formulas, with two
benchmark-specific conventions preserved: University-1652 follows its
official evaluator's distractor and junk-image handling and AP computation,
and DenseUAV additionally reports SDM@K — the distance-weighted metric it
introduced \cite{dai2024denseuav}, computed from the GPS coordinates of the
retrieved locations.

\section{Implementation Details}
The implementation is written in Python using PyTorch, \texttt{timm}, FAISS,
rasterio, matplotlib, OpenCV, Kornia, and PyQt5. CUDA and automatic mixed
precision are used when a CUDA-enabled PyTorch build is available. The
training pipeline parameterizes epochs, batch size, GPU micro-batch size,
learning rate, weight decay, device selection, worker count, mixed precision,
checkpoint resumption, data shuffling, backbone-tail unfreezing,
target cluster count, loss weights, backbone selection, and
adaptive cluster rebuilding.

The evaluation interface separates parameter configuration from result
inspection. It supports VisLoc sequence selection, automatic map and CSV
detection, checkpoint-aware feature caching, and
comparison runs between checkpoints. It visualizes the map cluster matrix as a
two-dimensional tile grid where cells with the same color belong to the same
embedding cluster.
Clusters with low descriptor variance can be highlighted as diagnostics,
allowing the user to inspect routing failure without removing those tiles from
retrieval. The interface also
supports a comparison model with an independent checkpoint and cluster count,
which allows controlled evaluation of alternative $K$ values or training runs
on the same map and query sequence.
Query-level inspection reports the original UAV image, the matched crop, the
top-$K$ matching crop when applicable, and the crop nearest to ground truth.

\section{Experimental Protocol}
The method is evaluated on five benchmarks, grouped into two protocol
families. The \emph{open-map} family applies the Game4Loc tiling and
matching protocol \cite{ji2024game4loc}: UAV-VisLoc
\cite{xu2024uavvisloc}, whose large georeferenced maps, pose metadata, and
diverse real terrain match the intended operating condition of the proposed
pipeline, and GTA-UAV, evaluated with Game4Loc's own metadata and evaluator
conventions. The \emph{closed-gallery} family follows each benchmark's
official evaluator — DenseUAV, SUES-200, and University-1652 — enabling
direct comparison against published state-of-the-art numbers
(Tables~\ref{tab:soa_sues200}--\ref{tab:soa_denseuav}). In every case the
evaluated model is the one trained on that benchmark's own official
training split (Section~\ref{sec:training_data}).
Table~\ref{tab:config} summarizes the shared evaluation configuration.

\subsection{Game4Loc Evaluation Protocol Applied to UAV-VisLoc}
The Game4Loc evaluation protocol is applied as follows.
The satellite gallery for each UAV-VisLoc area is the pre-tiled image pyramid
shipped with the dataset (tile naming: \texttt{\{area\}\_\{z\}\_\{x\}\_\{y\}.png},
$256\times256$ px, zoom levels $z_{\max}-3$ to $z_{\max}$).
For each drone query, the altitude and heading from the sequence CSV file are
used to compute the camera footprint with FOV$_H=36^\circ$ and FOV$_V=52^\circ$.
Positive tiles are those with IoU $\geq 0.39$ or tile recall $\geq 0.80$
against this footprint.
Retrieval is performed in drone-to-satellite (D2S) direction.
Same-area evaluation uses the sequences whose satellite maps are included in
the gallery; cross-area evaluation uses a held-out area.
Reported metrics are Recall@1, Recall@5, AP, SDM@1, SDM@5, Dis@1, Dis@5,
and query throughput (FPS). Table~\ref{tab:game4loc_condition} summarizes the
protocol.

The training condition differs from Game4Loc's own training data: whereas
Game4Loc trains on GTA-UAV synthetic drone-to-satellite pairs, the proposed
system trains on real UAV-VisLoc D2S pairs (using the same Game4Loc tiling
protocol) as described in Section~\ref{sec:training_data}. The evaluation
protocol, gallery structure, and metrics are identical to Game4Loc.

\subsection{GTA-UAV Protocol}
GTA-UAV itself is evaluated with Game4Loc's own evaluator conventions on
the same-area split: the gallery is the single flat set of all satellite
tiles across the four zoom levels, with no per-area or per-zoom masking;
queries are the test drone images with at least one positive tile; and only
the metadata's positive list (IoU $\geq 0.39$) counts as a hit, even though
training also uses semi-positives. Tile centre coordinates for the distance
metrics are recovered from the tile naming convention exactly as in the
official evaluator, and Recall@K, AP, SDM@K, and Dis@K follow the official
formulas.

\subsection{DenseUAV Protocol}
DenseUAV is evaluated on its official test split: 2,331 drone queries
(777 held-out locations at three altitudes) are scored against the full
18,198-tile satellite gallery in a single pass, with a hit defined as
retrieving a tile of the query's exact location ID. Recall@K, AP, and
DenseUAV's spatial distance metric SDM@K are reported.

\subsection{SUES-200 Protocol}
SUES-200 uses the official 120/80 train/test location split: 16,000 test
queries (80 locations $\times$ 4 altitudes $\times$ 50 images) are
evaluated per altitude against the 80 test locations' satellite tiles — the
gallery convention used by the published comparisons — and additionally
against the harder all-200-tile confusion gallery in which the training
locations' tiles act as distractors (Table~\ref{tab:sues200_detail}).

\subsection{University-1652 Protocol}
University-1652 is evaluated in both directions from a single set of
weights: drone$\to$satellite ($\sim$37,855 drone queries against the
951-image satellite gallery, distractors included) and
satellite$\to$drone (701 satellite queries against $\sim$51,355 drone
images). Distractor and junk handling and the AP computation follow the
official benchmark evaluator.

\begin{table}[htbp]
\caption{Game4Loc Protocol Applied to UAV-VisLoc}
\label{tab:game4loc_condition}
\centering
\footnotesize
\setlength{\tabcolsep}{3pt}
\begin{tabular}{>{\raggedright\arraybackslash}p{0.40\columnwidth}
                >{\raggedright\arraybackslash}p{0.52\columnwidth}}
\toprule
Component & Setting \\
\midrule
Dataset & UAV-VisLoc (real UAV imagery, 10 areas) \\
Reference gallery & Pre-tiled $256{\times}256$ px pyramid, 4 zoom levels per area \\
Positive criterion & IoU $\geq 0.39$ or tile recall $\geq 0.80$ \\
Semi-positive criterion & IoU $\geq 0.14$ \\
Camera model & FOV$_H=36^\circ$, FOV$_V=52^\circ$; altitude and heading from CSV \\
Retrieval direction & Drone-to-satellite (D2S) \\
Training data & UAV-VisLoc D2S pairs (Game4Loc tiling, pos + semi-pos IoU weights) \\
Metrics & Recall@1/5, AP, SDM@1/5, Dis@1/5, FPS \\
\bottomrule
\end{tabular}
\end{table}

\begin{table}[htbp]
\caption{Evaluation Configuration}
\label{tab:config}
\centering
\footnotesize
\setlength{\tabcolsep}{3pt}
\begin{tabular}{>{\raggedright\arraybackslash}p{0.46\columnwidth}
                >{\raggedright\arraybackslash}p{0.46\columnwidth}}
\toprule
Parameter & Value \\
\midrule
Backbone & Swin-B (primary) or ConvNeXt-T, last stages unfrozen \\
Retrieval descriptor & Raw L2-normalized backbone output: 1024D (Swin-B), 768D (ConvNeXt-T); no head \\
Default number of clusters & 8--16 (configurable) \\
Map tile gallery & Game4Loc pre-tiled: $256\times256$ px at $z_{\max}-3$ to $z_{\max}$ zoom levels \\
Positive IoU threshold & $\geq 0.39$ (semi-positive $\geq 0.14$; tile recall $\geq 0.80$) \\
Training loss & $1.0$ InfoNCE + $1.0$ variance + $0.2$ hard negative margin + $2.0$ consistency \\
Hit criterion & Retrieved tile has IoU $\geq 0.39$ with drone footprint \\
\bottomrule
\end{tabular}
\end{table}

\section{Results and Discussion}
The experimental comparison is designed to isolate three claims: cluster
routing reduces retrieval cost, cluster-constrained candidate scoring improves
within-cluster matching, and the resulting speed-accuracy trade-off generalizes
beyond a single UAV-VisLoc sequence. The reported metrics follow the schema
exported by \texttt{eval\_history\_runs.json}: Recall@1, Recall@K, AP,
SDM@1, SDM@K, distance error, query throughput, tile/stride setting, and
cluster count.

Table~\ref{tab:ablation} isolates the paper's central mechanism -- does
cluster-structured hard-negative mining actually help, holding everything
else fixed -- across all three backbones on DenseUAV, and separately tests
whether the learned descriptor and cluster routing improve over flat
descriptor search on GTA-UAV.

\begin{table}[htbp]
\caption{Ablation Study. DenseUAV clustering ablation, 3 backbones (top) and
GTA-UAV same-area under Game4Loc tiling (bottom). Bold = best per section
per metric. $\dagger$~SDM@3 (Game4Loc paper reports SDM@3; ours are SDM@1).}
\label{tab:ablation}
\centering
\footnotesize
\setlength{\tabcolsep}{3pt}
\begin{tabular}{>{\raggedright\arraybackslash}p{0.33\columnwidth}
                >{\centering\arraybackslash}p{0.11\columnwidth}
                >{\centering\arraybackslash}p{0.11\columnwidth}
                >{\centering\arraybackslash}p{0.11\columnwidth}
                >{\centering\arraybackslash}p{0.12\columnwidth}
                >{\centering\arraybackslash}p{0.12\columnwidth}}
\toprule
Variant & R@1 & R@5 & AP & SDM@1 & FPS \\
\midrule
\multicolumn{6}{l}{\textit{DenseUAV (drone$\to$satellite, same-area)}} \\
\midrule
% TODO: fill in from denseuav_ablation sweep (fresh ImageNet-start, 20ep,
% single variable = enable_clustering; see scratchpad/denseuav_ablation/)
Swin-B, no clustering & TBD & TBD & -- & TBD & -- \\
Swin-B, clustered (ours) & TBD & TBD & -- & TBD & -- \\
ConvNeXt-T, no clustering & TBD & TBD & -- & TBD & -- \\
ConvNeXt-T, clustered (ours) & TBD & TBD & -- & TBD & -- \\
ConvNeXt-B, no clustering & TBD & TBD & -- & TBD & -- \\
ConvNeXt-B, clustered (ours) & TBD & TBD & -- & TBD & -- \\
\midrule
\multicolumn{6}{l}{\textit{GTA-UAV same-area (Game4Loc tiling)}} \\
\midrule
Game4Loc~\cite{ji2024game4loc} (paper) & 0.8495 & 0.9759 & \textbf{0.9015} & $0.8803^\dagger$ & -- \\
Flat descriptor, cosine (ours) & 0.8874 & \textbf{0.9868} & -- & \textbf{0.8473} & -- \\
Two-stage, cosine (ours) & \textbf{0.9007} & 0.9834 & -- & 0.8340 & -- \\
\bottomrule
\end{tabular}
\end{table}

Table~\ref{tab:cluster_tradeoff} reports representative evaluation-history
rows showing how the routing cluster count and searched-cluster width relate
to recall and throughput. These rows are not a fully controlled hyperparameter
sweep, but they match the metrics emitted by the current evaluator and provide
the reporting format for a controlled sweep.

\begin{table}[htbp]
\caption{Cluster Granularity and Search-Width Trade-off}
\label{tab:cluster_tradeoff}
\centering
\footnotesize
\setlength{\tabcolsep}{3pt}
\begin{tabular}{>{\centering\arraybackslash}p{0.11\columnwidth}
                >{\centering\arraybackslash}p{0.13\columnwidth}
                >{\centering\arraybackslash}p{0.13\columnwidth}
                >{\centering\arraybackslash}p{0.13\columnwidth}
                >{\centering\arraybackslash}p{0.13\columnwidth}
                >{\centering\arraybackslash}p{0.13\columnwidth}
                >{\centering\arraybackslash}p{0.13\columnwidth}}
\toprule
Run & $K$ & Seq. & R@1 & R@5 & SDM@1 & FPS \\
\midrule
72 & 8 & 04 & 0.5325 & 0.9499 & 0.4132 & 251.3 \\
87 & 16 & 04 & 0.9187 & 1.0000 & 0.6895 & 499.5 \\
3 & 32 & 04 & 0.6314 & 0.9539 & -- & 401.3 \\
\bottomrule
\end{tabular}
\end{table}

Table~\ref{tab:cost_breakdown} separates one-time map preprocessing from
per-query retrieval. This distinction is necessary because map tiling,
embedding extraction, and FAISS clustering can be cached, whereas query
encoding and candidate scoring determine online localization latency.

\begin{table}[htbp]
\caption{Computational Cost Breakdown}
\label{tab:cost_breakdown}
\centering
\footnotesize
\setlength{\tabcolsep}{3pt}
\begin{tabular}{>{\raggedright\arraybackslash}p{0.42\columnwidth}
                >{\centering\arraybackslash}p{0.22\columnwidth}
                >{\centering\arraybackslash}p{0.24\columnwidth}}
\toprule
Component & Runtime role & Cache behavior \\
\midrule
GeoTIFF tiling & One-time preprocessing & Cached by map and tile geometry \\
Map descriptor embedding & One-time preprocessing & Cached by checkpoint hash \\
FAISS map K-means & One-time preprocessing & Cached with tile assignments \\
Query embedding & Per-query online & Cached by checkpoint and image \\
Candidate scoring & Per-query online & Not cached across new queries \\
\bottomrule
\end{tabular}
\end{table}

Table~\ref{tab:sequence_generalization} evaluates whether the method remains
effective on UAV-VisLoc sequences not used for method tuning. If similar
datasets such as DenseUAV, SUES-200, or University-1652 are used, they can be
reported in the same format with dataset-specific retrieval metrics.

\begin{table}[htbp]
\caption{Cross-Sequence Generalization}
\label{tab:sequence_generalization}
\centering
\footnotesize
\setlength{\tabcolsep}{3pt}
\begin{tabular}{>{\centering\arraybackslash}p{0.10\columnwidth}
                >{\centering\arraybackslash}p{0.12\columnwidth}
                >{\centering\arraybackslash}p{0.12\columnwidth}
                >{\centering\arraybackslash}p{0.12\columnwidth}
                >{\centering\arraybackslash}p{0.12\columnwidth}
                >{\centering\arraybackslash}p{0.13\columnwidth}
                >{\centering\arraybackslash}p{0.13\columnwidth}}
\toprule
Seq. & Queries & R@1 & R@5 & AP & Median hit & Tile/stride \\
\midrule
03 & 768 & 0.7747 & 1.0000 & 0.5932 & 110.0 m & 340/110 \\
04 & 738 & 0.9187 & 1.0000 & 0.7188 & 129.9 m & 400/130 \\
05 & 473 & 0.4397 & 0.9535 & -- & -- & 260/90 \\
07 & 28 & 0.4286 & 0.7500 & 0.6118 & 50.1 m & 140/50 \\
11 & 590 & 0.6034 & 0.9966 & 0.4468 & 190.0 m & 570/190 \\
\bottomrule
\end{tabular}
\end{table}

% ── SUES-200 state-of-the-art comparison ─────────────────────────────────────
\begin{table*}[htbp]
\caption{Comparison with state-of-the-art on SUES-200 (Recall@1, \%, drone$\to$satellite).
         SOTA values as reported in~\cite{li2025cvd}, except
         MEAN~\cite{chen2025mean} taken from its own paper.
         Ours: official 120/80 train/test location split, evaluated on the 80
         held-out test locations against their 80 satellite tiles — the same
         gallery convention used by the compared methods'
         evaluations\textsuperscript{\S}.}
\label{tab:soa_sues200}
\centering
\footnotesize
\setlength{\tabcolsep}{8pt}
\begin{tabular}{l l cccc}
\toprule
Method & Backbone & 150\,m & 200\,m & 250\,m & 300\,m \\
\midrule
LPN~\cite{wang2021lpn}                 & ResNet-50   & 58.20 & 69.60 & 75.60 & 78.50 \\
FSRA~\cite{dai2022fsra}                & ViT-S       & 78.70 & 85.65 & 88.95 & 91.50 \\
MCCG~\cite{shen2024mccg}               & ConvNeXt-T  & 79.96 & 85.01 & 90.47 & 94.20 \\
Sample4Geo~\cite{deuser2023sample4geo} & ConvNeXt-B  & 94.75 & 96.75 & 97.25 & 97.20 \\
Game4Loc~\cite{ji2024game4loc}        & ConvNeXt-B  & 94.62 & 96.55 & 97.55 & 97.67 \\
MEAN~\cite{chen2025mean}               & ConvNeXt-T  & 95.50 & \textbf{98.38} & 98.95 & 99.17 \\
DAC~\cite{xia2024dac}\textsuperscript{$\ddagger$}                & ConvNeXt-B  & \textbf{96.80} & 97.48 & 98.20 & 97.58 \\
\midrule
\textbf{Ours}\textsuperscript{\S} & Swin-B & 96.30 & 98.25 & \textbf{99.00} & 99.15 \\
\textbf{Ours}\textsuperscript{\S,\dag} & ConvNeXt-T & 91.77 & 95.20 & 96.90 & 97.52 \\
\textbf{Ours}\textsuperscript{\S,\dag} & ConvNeXt-B & 94.97 & 97.28 & 98.50 & \textbf{99.35} \\
\bottomrule
\end{tabular}

\smallskip
\raggedright\footnotesize
\textsuperscript{\S}\,Under the harder all-200-tile confusion gallery (the 120
training locations' tiles added as distractors), Swin-B scores 91.72 / 94.15 /
95.85 / 97.02 (overall R@1 94.69, mAP 96.97), ConvNeXt-T scores 87.10 /
91.72 / 94.10 / 95.58 (overall R@1 92.12, mAP 95.08), and ConvNeXt-B scores
93.00 / 95.53 / 97.02 / 97.82 (overall R@1 95.84, mAP 97.46); the compared
methods do not report this setting. Table~\ref{tab:sues200_detail} reports
the full per-altitude results for all models and gallery settings.\\
\textsuperscript{\dag}\,ConvNeXt rows: the DenseUAV-trained epoch-25
checkpoints (Table~\ref{tab:soa_denseuav}) fine-tuned on SUES-200's own
official training split (lr=$1\times10^{-5}$, last 2 backbone stages
unfrozen, 20 epochs) -- the same fine-tuning recipe used for the Swin-B row.\\
\textsuperscript{$\ddagger$}\,DAC's SUES-200 numbers are taken from
MEAN's~\cite{chen2025mean} comparison table; DAC's own paper was not
accessed directly (paywalled, no preprint found) and its official
repository's README includes no results table, though it does release
trained checkpoints, unlike CEUSP.
\end{table*}

% ── SUES-200 detailed per-altitude results (ours) ─────────────────────────────
\begin{table*}[htbp]
\caption{Detailed SUES-200 results (\%, drone$\to$satellite, official
         120/80 split, 16\,000 test queries) for all three of our backbones.
         Left: the 80-tile test gallery used by the compared methods in
         Table~\ref{tab:soa_sues200}. Right: the all-200-tile confusion
         gallery, in which the 120 training locations' satellite tiles act as
         distractors — a harder setting closer to open-map operation.
         Recall degrades toward lower altitudes in both settings (smaller
         ground footprint, less layout context per query) for every backbone,
         and the confusion gallery costs $2.0$--$4.7$ R@1 points at 150\,m but
         only $1.4$--$2.1$ at 300\,m.}
\label{tab:sues200_detail}
\centering
\footnotesize
\setlength{\tabcolsep}{7pt}
\begin{tabular}{l ccc ccc}
\toprule
& \multicolumn{3}{c}{80-tile test gallery}
& \multicolumn{3}{c}{200-tile confusion gallery} \\
\cmidrule(lr){2-4}\cmidrule(lr){5-7}
Altitude & R@1 & R@5 & R@10 & R@1 & R@5 & R@10 \\
\midrule
\multicolumn{7}{l}{\textit{Swin-B}} \\
150\,m   & 96.30 & 99.48  & 100.00 & 91.72 & 98.88  & 99.67 \\
200\,m   & 98.25 & 99.90  & 100.00 & 94.15 & 99.70  & 99.98 \\
250\,m   & 99.00 & 100.00 & 100.00 & 95.85 & 99.98  & 100.00 \\
300\,m   & 99.15 & 100.00 & 100.00 & 97.02 & 100.00 & 100.00 \\
\midrule
Overall  & 98.17 & 99.84  & 100.00 & 94.69 & 99.64  & 99.91 \\
mAP      & \multicolumn{3}{c}{98.90} & \multicolumn{3}{c}{96.97} \\
\midrule
\multicolumn{7}{l}{\textit{ConvNeXt-T}} \\
150\,m   & 91.77 & 98.47  & 99.17  & 87.10 & 97.38  & 98.78 \\
200\,m   & 95.20 & 99.28  & 99.60  & 91.72 & 98.92  & 99.42 \\
250\,m   & 96.90 & 99.45  & 99.72  & 94.10 & 99.22  & 99.60 \\
300\,m   & 97.52 & 99.65  & 99.88  & 95.58 & 99.33  & 99.80 \\
\midrule
Overall  & 95.35 & 99.21  & 99.59  & 92.12 & 98.71  & 99.40 \\
mAP      & \multicolumn{3}{c}{97.04} & \multicolumn{3}{c}{95.08} \\
\midrule
\multicolumn{7}{l}{\textit{ConvNeXt-B}} \\
150\,m   & 94.97 & 98.92  & 99.45  & 93.00 & 98.55  & 99.25 \\
200\,m   & 97.28 & 99.55  & 99.78  & 95.53 & 99.33  & 99.72 \\
250\,m   & 98.50 & 99.92  & 100.00 & 97.02 & 99.88  & 99.95 \\
300\,m   & 99.35 & 100.00 & 100.00 & 97.82 & 100.00 & 100.00 \\
\midrule
Overall  & 97.52 & 99.60  & 99.81  & 95.84 & 99.44  & 99.73 \\
mAP      & \multicolumn{3}{c}{98.43} & \multicolumn{3}{c}{97.46} \\
\bottomrule
\end{tabular}
\end{table*}

% ── University-1652 state-of-the-art comparison ───────────────────────────────
\begin{table*}[htbp]
\caption{Comparison with state-of-the-art on University-1652 (R@1 and AP, \%).
         D$\!\to\!$S: drone query, satellite gallery ($\sim$37\,855 queries, 951 gallery).
         S$\!\to\!$D: satellite query, drone gallery (701 queries,
         $\sim$51\,355 gallery).
         SOTA values from the official leaderboard~\cite{zheng2020university1652},
         except MEAN~\cite{chen2025mean} taken from its own paper and
         DAC~\cite{xia2024dac}\textsuperscript{$\ddagger$} taken from MEAN's
         comparison table (see Table~\ref{tab:soa_sues200} footnote for why).
         Ours: per backbone, one model fine-tuned on the official 701-class
         training split (the ConvNeXt rows start from their DenseUAV-trained
         checkpoints of Table~\ref{tab:soa_denseuav}, fine-tuned with
         lr=$1\times10^{-5}$, last 2 backbone stages unfrozen, 20 epochs --
         the same recipe as the Swin-B row), both directions evaluated from
         the same weights; distractor handling and the AP formula follow the
         official benchmark evaluator.}
\label{tab:soa_u1652}
\centering
\footnotesize
\setlength{\tabcolsep}{8pt}
\begin{tabular}{l l cc cc}
\toprule
& & \multicolumn{2}{c}{Drone $\to$ Satellite}
& \multicolumn{2}{c}{Satellite $\to$ Drone} \\
\cmidrule(lr){3-4}\cmidrule(lr){5-6}
Method & Backbone & R@1 & AP & R@1 & AP \\
\midrule
LPN~\cite{wang2021lpn}                 & ResNet-50  & 75.93 & 79.14 & 86.45 & 74.79 \\
FSRA~\cite{dai2022fsra}                & ViT-S      & 84.51 & 86.71 & 88.45 & 83.37 \\
MCCG~\cite{shen2024mccg}               & ConvNeXt-T & 89.64 & 91.32 & 94.30 & 89.39 \\
Sample4Geo~\cite{deuser2023sample4geo} & ConvNeXt-B & 92.65 & 93.81 & 95.14 & 91.39 \\
Game4Loc~\cite{ji2024game4loc}         & ConvNeXt-B & 91.32 & 92.56 & 94.43 & 90.83 \\
MEAN~\cite{chen2025mean}               & ConvNeXt-T & 93.55 & 94.53 & 96.01 & 92.08 \\
DAC~\cite{xia2024dac}\textsuperscript{$\ddagger$}                & ConvNeXt-B & 94.67 & \textbf{95.50} & 96.43 & 93.79 \\
\midrule
\textbf{Ours}\textsuperscript{$\P$} & Swin-B & \textbf{94.73} & 95.44 & \textbf{96.86} & \textbf{93.83} \\
\textbf{Ours} & ConvNeXt-T & 82.08 & 84.69 & 92.44 & 80.38 \\
\textbf{Ours} & ConvNeXt-B & 90.91 & 92.44 & 95.72 & 90.32 \\
\bottomrule
\end{tabular}

\smallskip
\raggedright\footnotesize
\textsuperscript{$\P$}\,R@1 from a cluster-count sweep winner (K=64,
resumed from the previously-reported checkpoint with the content-hash MES
fix active, 20 epoch-equivalents total). AP/S$\to$D columns are not yet
recomputed for this checkpoint and still show the prior (K=16)
checkpoint's values.
\end{table*}

% ── DenseUAV state-of-the-art comparison ──────────────────────────────────────
\begin{table}[htbp]
\caption{Comparison with state-of-the-art on DenseUAV (drone$\to$satellite).
         SDM@1 is the DenseUAV spatial-distance metric.
         Comparison values as reported in~\cite{xu2025ceusp} and the
         original papers, except MCCG~\cite{shen2024mccg}\textsuperscript{$\triangle$}.}
\label{tab:soa_denseuav}
\centering
\footnotesize
\setlength{\tabcolsep}{3pt}
\resizebox{\columnwidth}{!}{%
\begin{tabular}{l l c ccc}
\toprule
Method & Backbone & Params & R@1 & R@5 & SDM@1 \\
\midrule
DenseUAV baseline~\cite{dai2024denseuav}    & ViT-S      & $\sim$22\,M & 83.01 & 95.58 & 86.50 \\
Sample4Geo~\cite{deuser2023sample4geo}      & ConvNeXt-B & 87.6\,M     & 49.38 & 78.29 & 61.72 \\
MCCG~\cite{shen2024mccg}\textsuperscript{$\ast$}             & ConvNeXt-T & $\sim$28\,M & 83.14 & --    & --    \\
MCCG~\cite{shen2024mccg}\textsuperscript{$\triangle$}        & ConvNeXt-T & $\sim$28\,M & 90.26 & \textbf{97.90} & \textbf{91.82} \\
CAMP~\cite{wu2024camp}\textsuperscript{$\diamond$}                        & ConvNeXt-B & 91.4\,M     & 88.72 & --    & --    \\
CEUSP~\cite{xu2025ceusp}\textsuperscript{$\star$}         & ConvNeXt-T & $\sim$28\,M & 89.45 & 96.05 & 91.01 \\
DINOv2-GLFA+CESP~\cite{yang2025dinov2}      & DINOv2-B   & $\sim$86\,M & 86.27 & 96.83 & 88.87 \\
\midrule
\textbf{Ours}\textsuperscript{\S}                & Swin-B & $\sim$88\,M & 86.62 & 96.83 & 89.50 \\
\textbf{Ours + VQ re-rank}\textsuperscript{\S\P} & Swin-B & $\sim$88\,M & 88.55 & 96.61 & --\textsuperscript{\dag} \\
\textbf{Ours}\textsuperscript{\S\ddag}                & ConvNeXt-T & $\sim$28\,M & 86.14 & 96.91 & 89.12 \\
\textbf{Ours + VQ re-rank}\textsuperscript{\S\P\ddag} & ConvNeXt-T & $\sim$28\,M & \textbf{94.38} & \textbf{98.24} & --\textsuperscript{\dag} \\
\textbf{Ours}\textsuperscript{\S\#}                & ConvNeXt-B & $\sim$89\,M & 87.22 & 96.95 & 89.85 \\
\textbf{Ours + VQ re-rank}\textsuperscript{\S\P\#} & ConvNeXt-B & $\sim$89\,M & 92.58 & 98.16 & --\textsuperscript{\dag} \\
\bottomrule
\end{tabular}}

\vspace{2pt}
{\footnotesize
\textsuperscript{\ddag}\,Selected ConvNeXt-T checkpoint: continued from this
paper's originally-reported epoch-25 DenseUAV checkpoint for 5 more epochs
at lr=$2\times10^{-5}$ with the last 2 backbone stages unfrozen (winner of a
3-config unfreeze-depth sweep at this learning rate over a 3-stage-unfrozen
alternative, which scored 85.50 R@1; a 4-stage-unfrozen configuration was
not run). Also reports R@10\,$=98.71$, mAP\,$=74.18$ (coarse) and
R@10\,$=98.71$, mAP\,$=79.42$ (VQ re-rank).\\
The coarse and VQ re-rank numbers for both ConvNeXt rows in this table were
revised upward from an earlier draft (ConvNeXt-T: 82.41/91.89; ConvNeXt-B:
84.56/91.29) after fixing a bug in \texttt{model.py}'s shared
retrieval-feature helper used by the VQ pipeline
(\texttt{encode\_spatial\_features}/\texttt{encode\_features\_and\_scale}):
for ConvNeXt-style backbones, timm applies a final LayerNorm2d in the
classifier head after global pooling that this helper skipped, producing a
materially different, weaker embedding than the model's actual retrieval
output (mean absolute difference $\approx0.75$ against the correct
embedding). Swin-B, ViT, and DINOv2 backbones are unaffected -- verified
their \texttt{forward\_features} output already matches the full backbone
forward exactly.\\
\textsuperscript{$\star$}\,CEUSP reports results at $256\times256$ input
(ours: $384\times384$); at the time of writing, no public code or pretrained
weights were available for CEUSP, so its numbers could not be independently
verified.\\
\textsuperscript{$\ast$}\,MCCG's DenseUAV R@1 as reported secondhand via
CEUSP's comparison table \cite{xu2025ceusp} -- MCCG's own paper does not
evaluate DenseUAV, and CEUSP states no reproduction methodology (epochs,
learning rate, or whether official code was used) for this or any other
baseline it reports, nor R@5/SDM@1 for MCCG. Superseded by our own
independently verified reproduction on the row below.\\
\textsuperscript{$\triangle$}\,Our own reproduction, trained from MCCG's
official released code \cite{shen2024mccg} for 199 epochs with the paper's
stated recipe (lr=$0.01$, batch size 8, triplet-loss weight 0.3), evaluated
under the same protocol as every other row in this table.\\
\textsuperscript{$\diamond$}\,CAMP's DenseUAV figure is taken from CEUSP's
comparison table~\cite{xu2025ceusp}; CAMP's own paper evaluates
University-1652 and SUES-200 only and reports no DenseUAV results, so R@5
and SDM@1 are unavailable.\\
\textsuperscript{\S}\,Same-area retrieval, exact location-ID hit over the full
18{,}198-tile confusion gallery (single pass) — the same hit criterion used for
the other rows in this table. Also reports R@10\,$=98.63$ and mAP\,$=77.25$.\\
\textsuperscript{\P}\,With the vector-quantized semantic re-ranking stage of
Section~\ref{sec:vq_rerank} applied to the coarse top-10 (also reports
R@10\,$=98.63$ and mAP\,$=78.69$).\\
\textsuperscript{\dag}\,SDM (spatial-distance metric) not separately computed
for the re-ranked ordering.\\
\textsuperscript{\#}\,Selected ConvNeXt-B checkpoint (epoch 25 of a chunked
training chain from ImageNet, transplanted recipe lr=$2\times10^{-4}$, last 3
backbone stages unfrozen); per-epoch selection used a lighter-weight
monitoring eval (R@1 76.53 / 80.44 / 81.25 / 85.16 / 87.22 at epochs
5/10/15/20/25, epoch 30 regressed to 86.83) to pick the checkpoint cheaply --
this table's row now matches that monitoring trajectory exactly (87.22)
after the retrieval-feature bug fix described under \ddag; the same
checkpoint's earlier-draft number (84.56) predates the fix. Also reports
R@10\,$=98.84$, mAP\,$=76.95$ (coarse) and R@10\,$=98.84$, mAP\,$=80.75$
(VQ re-rank).}
\end{table}

The most important training diagnostic is the anchor-positive same-cluster
probability. This value measures whether the descriptor places paired
cross-source crops into the same K-means region. The cluster-consistency loss,
variance loss, and same-cluster negative similarity provide direct feedback on
whether the routing descriptor is both stable and discriminative.

\section{Limitations}
\label{sec:limitations}
The training pairs are generated from co-registered GeoTIFF crops with shared
ground geometry and therefore do not fully model oblique UAV-to-satellite
viewpoint changes. The frozen backbone
reduces training cost and stabilizes descriptors but may limit adaptation to
domain-specific imagery. The K-means clustering step may remain
expensive for very large maps, although tile images, embeddings, and cluster
data are cached. Finally, evaluation uses an IoU $\geq 0.39$ hit criterion
following Game4Loc; future work could complement this with stricter
localization accuracy thresholds such as sub-10 m error rates.

\section{Conclusion}
This paper presented a cluster-routed descriptor method for UAV
cross-view localization. The raw L2-normalized output of a pretrained
backbone (Swin-B, or ConvNeXt-T for a $3\times$ smaller model), fine-tuned
through its last stages, serves as a single descriptor for learned routing,
candidate scoring, and — centrally — cluster-structured hard-negative mining:
negatives drawn from a query's own and nearest clusters, with co-positive
exclusion, concentrate the contrastive signal on visually repetitive
look-alike regions. The training pipeline
combines paired InfoNCE, cluster-consistency learning, variance
regularization, and same-cluster hard-negative suppression, while the
evaluation pipeline performs scalable
coarse-to-fine search over a Game4Loc multi-zoom tile gallery with an
optional vector-quantized semantic re-ranking stage. The method provides
an operational foundation for UAV localization experiments with explicit
cluster monitoring, checkpoint-aware cached evaluation,
and query-level qualitative analysis.

\section*{Funding}
% TODO: replace with actual funding source(s), or keep as-is if none.
This research did not receive any specific grant from funding agencies in
the public, commercial, or not-for-profit sectors.

\section*{CRediT Author Statement}
% TODO: assign roles per author using the CRediT taxonomy
% (Conceptualization, Methodology, Software, Validation, Formal analysis,
% Investigation, Resources, Data curation, Writing - original draft,
% Writing - review \& editing, Visualization, Supervision, Project
% administration, Funding acquisition). Example below -- edit per author.
\textbf{Kim-Phuong Phung}: Conceptualization, Methodology, Software,
Validation, Formal analysis, Investigation, Data curation, Writing -
original draft, Visualization. \textbf{Quang-Uy Nguyen}: Conceptualization,
Writing - review \& editing, Supervision, Project administration.

\section*{Declaration of Generative AI Use}
During the preparation of this work, the author(s) used Claude Code
(Anthropic) to assist with code implementation, experiment execution, and
manuscript writing. All AI-assisted code, experimental results, and text
were carefully reviewed, edited, and verified by the author(s), who take
full responsibility for the content of this article.

\section*{Data Availability Statement}
Code is available at
\url{https://github.com/kimphg/UAV_Visual_Clustering}. Trained model
checkpoints are available at
\url{https://drive.google.com/drive/folders/1MJjMTtpWYBR8V7dTb40Exr1q-efToXzD}.
Datasets used (DenseUAV, SUES-200, University-1652, GTA-UAV) are publicly
available from their respective original sources cited in this paper.

\bibliographystyle{IEEEtran}
\bibliography{references}

\end{document}

